% !TEX root = ../高等概率论.tex
% =====================================================================
%  集合系：代数、σ-代数、生成的 σ-代数、Borel σ-代数、π-类与 λ-类、π-λ 定理
% =====================================================================
\lecture{集合运算与示性函数，代数与 \texorpdfstring{$\sigma$}{σ}-代数，生成的 \texorpdfstring{$\sigma$}{σ}-代数，Borel \texorpdfstring{$\sigma$}{σ}-代数，\texorpdfstring{$\pi$}{π}-类与 \texorpdfstring{$\lambda$}{λ}-类，Dynkin 的 \texorpdfstring{$\pi$-$\lambda$}{π-λ} 定理}
\label{lec:set-systems}

我们先描述一下这门课所要建立的理论的大体轮廓。初等概率论中处理的大多是有限的模型：掷骰子、抽签、有限次抛硬币。样本空间只有有限个点，每一个子集都可以谈论概率，而计算概率归根到底是数数。可是一旦样本空间变成无限的，比如“在 $[0,1]$ 中随机地取一个数”，或者“把一枚硬币无穷次地抛下去”，这套做法就行不通了：单独一个点的概率只能是 $0$，概率不能再通过“把每个点的概率加起来”得到。我们需要一种新的语言，这就是测度论。

按照 Kolmogorov 的公理化，概率论的数学模型是一个三元组 $(\Omega,\cF,\P)$，称为\term{概率空间}（probability space）。其中 $\Omega$ 是全部可能结果组成的集合；$\cF$ 是 $\Omega$ 的某些子集组成的集合，我们只对 $\cF$ 中的集合谈论概率；$\P$ 是定义在 $\cF$ 上的函数，给每个集合指定它的概率。三样东西各自回答一个问题：可能发生什么，哪些事情可以问概率，概率是多少。这一讲只讨论第二个问题，即 $\cF$ 应当满足什么条件。

一个自然的疑问是：为什么不干脆取 $\cF=2^\Omega$，允许对一切子集谈论概率？答案是做不到。以 $[0,1)$ 上的“均匀分布”为例，我们希望区间的概率等于它的长度，并且希望概率在平移下不变、对可数个不交集合可加。第~\ref{lec:extension-of-measures} 讲的作业将证明，满足这些要求的函数不可能定义在 $[0,1)$ 的全体子集上。所以我们只能退一步，只在一部分“足够好”的集合上定义概率。另一方面，这一部分集合又不能太少：概率论中关心的许多事件，比如“抛硬币时正面出现无穷多次”“频率 $S_n/n$ 收敛到 $1/2$”，都要通过对可数个事件取交、取并才能写出来。这两方面的要求折中的结果，就是本讲的 $\sigma$-代数。

我们用一个具体的问题贯穿这一讲。把一枚硬币无穷次地抛下去，记 $\omega_n=1$ 表示第 $n$ 次是正面，$\omega_n=0$ 表示反面，一次试验的结果就是一个无穷的 $0$-$1$ 序列 $\omega=(\omega_1,\omega_2,\dots)$，样本空间是
\[
\Omega=\{0,1\}^{\N}=\{\omega=(\omega_1,\omega_2,\dots):\ \omega_n\in\{0,1\}\}.
\]
记 $S_n(\omega)=\omega_1+\cdots+\omega_n$ 为前 $n$ 次中正面的次数。强大数定律要说的是，事件
\[
\Bigl\{\omega:\ \lim_{n\to\infty}\frac{S_n(\omega)}n=\frac12\Bigr\}
\]
的概率等于 $1$。在谈论它的概率之前，我们至少要知道它是一个“可以谈论概率的集合”。这件事并不是显然的：“第 $n$ 次是正面”只涉及一次抛掷，而上面这个集合涉及全部无穷多次抛掷。本讲的例题（第~\pageref{we:set-systems:slln-event} 页）将说明，它可以由“第 $n$ 次是正面”这样的简单事件经过可数次取并、取交、取补得到。这正是 $\sigma$-代数要对可数次运算封闭的原因。

这一讲的内容几乎全是集合论的语言，初读时可能有些枯燥。但后面几乎每一个定理都要用到这里的工具，特别是最后一节的 $\pi$-$\lambda$ 定理，请同学们务必把它的证明搞清楚。

\topic{本讲用到的前文结果}

这是第一讲，不用到前文的结果。我们只用到数学分析中下面几条事实，读者如果忘了，可以先在这里复习一下。
\begin{enumerate}
\item \emph{可数集。}称集合是\term{可数的}（countable），如果它是有限集，或者能与 $\N$ 一一对应。可数集的子集可数；可数个可数集的并可数；$\Q$、$\Q\times\Q$、$\Q^d$ 都可数。$\R$ 与 $[0,1]$ 不可数。
\item \emph{有理数的稠密性。}任意两个实数 $a<b$ 之间都有有理数。
\item \emph{开集。}$U\subseteq\R$ 称为开集，若对每个 $x\in U$，存在 $\varepsilon>0$ 使 $(x-\varepsilon,x+\varepsilon)\subseteq U$。$\R^d$ 中把区间换成球 $\{y:\abs{y-x}<\varepsilon\}$。闭集是开集的补集。
\item \emph{数列的上极限与下极限。}$\limsup_na_n=\inf_n\sup_{m\geqslant n}a_m$，$\liminf_na_n=\sup_n\inf_{m\geqslant n}a_m$；数列收敛当且仅当二者相等且有限。
\end{enumerate}

\topic{集合运算与示性函数}

设 $\Omega$ 是一个非空集合，称为\term{样本空间}（sample space），其元素 $\omega$ 称为\term{样本点}（sample point）。$\Omega$ 的全体子集记作 $2^\Omega$，称为 $\Omega$ 的\term{幂集}（power set）。对 $A,B\subseteq\Omega$，记
\[
A^c=\Omega\setminus A,\qquad A\setminus B=A\cap B^c,\qquad A\triangle B=(A\setminus B)\cup(B\setminus A),
\]
分别称为 $A$ 的\term{补集}（complement）、$A$ 与 $B$ 的\term{差集}（difference）和\term{对称差}（symmetric difference）。若 $A\cap B=\varnothing$，称 $A,B$ \term{不交}（disjoint）；一列集合两两不交时，它们的并称为\term{不交并}（disjoint union），有时写作 $\sum_n A_n$ 或 $A_1+A_2+\cdots$。

对任意指标集 $I$ 和集合族 $\{A_i\}_{i\in I}$，\term{De Morgan 律}（De Morgan's laws）成立：
\[
\Bigl(\bigcup_{i\in I}A_i\Bigr)^c=\bigcap_{i\in I}A_i^c,\qquad \Bigl(\bigcap_{i\in I}A_i\Bigr)^c=\bigcup_{i\in I}A_i^c .
\]
事实上，$\omega\notin\bigcup_iA_i$ 当且仅当对每个 $i$ 都有 $\omega\notin A_i$，即 $\omega\in\bigcap_iA_i^c$；第二式对第一式两边取补并把 $A_i$ 换成 $A_i^c$ 即得。

\begin{definition}[示性函数]\label{def:set-systems:indicator}
设 $A\subseteq\Omega$。函数
\[
\one_A(\omega)=\begin{cases}1,&\omega\in A,\\0,&\omega\notin A\end{cases}
\]
称为 $A$ 的\term{示性函数}（indicator function）。
\end{definition}

示性函数把集合运算翻译成函数运算：
\[
\one_{A\cap B}=\one_A\one_B=\min\{\one_A,\one_B\},\qquad \one_{A\cup B}=\max\{\one_A,\one_B\},\qquad \one_{A^c}=1-\one_A;
\]
$A\subseteq B$ 当且仅当 $\one_A\leqslant\one_B$；$A_1,A_2,\dots$ 两两不交时 $\one_{\bigcup_nA_n}=\sum_n\one_{A_n}$。以后我们会经常在集合与函数之间来回转换：集合的问题有时写成示性函数更容易处理，积分的理论也正是从示性函数出发的。

\begin{example}[用示性函数验证集合等式]
我们验证 $\one_{A\cup B}=\one_A+\one_B-\one_A\one_B$。两边都只取 $0,1$ 两个值，逐点比较即可：$\omega$ 属于 $A,B$ 中恰好一个时，右边是 $1+0-0=1$；属于两个时是 $1+1-1=1$；都不属于时是 $0$。左边在三种情形下依次是 $1,1,0$。以后对这个等式两边取期望，就得到初等概率论中的加法公式 $\P(A\cup B)=\P(A)+\P(B)-\P(A\cap B)$。

再比如对称差：$\one_{A\triangle B}=\abs{\one_A-\one_B}=(\one_A-\one_B)^2$，这因为 $\omega\in A\triangle B$ 当且仅当 $\one_A(\omega)$ 与 $\one_B(\omega)$ 不同。于是 $A=B$ 当且仅当 $\one_A=\one_B$，“两个集合接近”可以用“两个函数接近”来刻画，以后讲测度的逼近时会用到这种想法。
\end{example}

\begin{definition}[集合列的上极限与下极限]\label{def:set-systems:limsup}
设 $A_1,A_2,\dots\subseteq\Omega$。定义
\[
\limsup_{n\to\infty}A_n=\bigcap_{n=1}^\infty\bigcup_{m=n}^\infty A_m,\qquad
\liminf_{n\to\infty}A_n=\bigcup_{n=1}^\infty\bigcap_{m=n}^\infty A_m,
\]
分别称为集合列 $\{A_n\}$ 的\term{上极限}（limit superior）与\term{下极限}（limit inferior）。若二者相等，称 $\{A_n\}$ 的\term{极限}存在，记作 $\lim_nA_n$。
\end{definition}

\begin{proposition}\label{prop:set-systems:limsup}
设 $A_1,A_2,\dots\subseteq\Omega$。则
\begin{enumerate}
\item $\omega\in\limsup_nA_n$ 当且仅当有无穷多个 $n$ 使 $\omega\in A_n$；因此概率论中常把 $\limsup_nA_n$ 写作 $\{A_n\ \text{i.o.}\}$（infinitely often，无穷多次发生）。
\item $\omega\in\liminf_nA_n$ 当且仅当存在 $N$，使得对一切 $n\geqslant N$ 都有 $\omega\in A_n$（$A_n$ 最终总发生）。
\item $\liminf_nA_n\subseteq\limsup_nA_n$，且 $(\limsup_nA_n)^c=\liminf_nA_n^c$。
\item $\one_{\limsup_nA_n}=\limsup_n\one_{A_n}$，$\one_{\liminf_nA_n}=\liminf_n\one_{A_n}$。
\item 若 $A_1\subseteq A_2\subseteq\cdots$，则 $\lim_nA_n=\bigcup_nA_n$；若 $A_1\supseteq A_2\supseteq\cdots$，则 $\lim_nA_n=\bigcap_nA_n$。
\end{enumerate}
\end{proposition}
\begin{proof}
(1) $\omega\in\bigcap_n\bigcup_{m\geqslant n}A_m$ 意味着：对每个 $n$，都存在 $m\geqslant n$ 使 $\omega\in A_m$。若只有有限多个 $m$ 使 $\omega\in A_m$，设其中最大的为 $M$（没有这样的 $m$ 时令 $M=0$），则取 $n=M+1$ 时找不到 $m\geqslant n$ 使 $\omega\in A_m$，矛盾。反之，若有无穷多个 $m$ 使 $\omega\in A_m$，则对每个 $n$，这些 $m$ 中必有不小于 $n$ 的，于是 $\omega\in\bigcup_{m\geqslant n}A_m$ 对一切 $n$ 成立。

(2) $\omega\in\bigcup_n\bigcap_{m\geqslant n}A_m$ 意味着存在 $n$，使 $\omega\in A_m$ 对一切 $m\geqslant n$ 成立，这正是所说的条件。

(3) 若 $\omega\in\liminf_nA_n$，由 (2) 存在 $N$ 使 $n\geqslant N$ 时 $\omega\in A_n$，这当然是无穷多个 $n$，由 (1) 得 $\omega\in\limsup_nA_n$。第二个等式由 De Morgan 律两次得到：
\[
\Bigl(\bigcap_n\bigcup_{m\geqslant n}A_m\Bigr)^c=\bigcup_n\Bigl(\bigcup_{m\geqslant n}A_m\Bigr)^c=\bigcup_n\bigcap_{m\geqslant n}A_m^c .
\]

(4) 数列 $\one_{A_n}(\omega)$ 只取 $0,1$ 两个值。它的上极限等于 $1$ 当且仅当它有无穷多项等于 $1$，即 $\omega$ 属于无穷多个 $A_n$，由 (1) 即 $\omega\in\limsup_nA_n$；否则上极限为 $0$。下极限等于 $1$ 当且仅当从某项起全等于 $1$，由 (2) 即 $\omega\in\liminf_nA_n$。

(5) 设 $A_n$ 递增。对每个 $n$，$\bigcap_{m\geqslant n}A_m=A_n$，$\bigcup_{m\geqslant n}A_m=\bigcup_{m\geqslant 1}A_m$，于是 $\liminf_nA_n=\bigcup_nA_n=\limsup_nA_n$。递减的情形类似：$\bigcap_{m\geqslant n}A_m=\bigcap_{m\geqslant1}A_m$，$\bigcup_{m\geqslant n}A_m=A_n$，两者都给出 $\bigcap_nA_n$。
\end{proof}

\begin{remark}[上极限与下极限的概率意义]
我们在概率论中关心的很多事件恰好是这种形式。假设 $A_n$ 是“第 $n$ 次抛硬币得到正面”这一事件，那么
\begin{enumerate}
\item $\limsup_nA_n$ 是“正面出现无穷多次”；
\item $\liminf_nA_n$ 是“从某一次起以后全是正面”；
\item $(\limsup_nA_n)^c=\liminf_nA_n^c$ 说的是：“正面只出现有限次”与“从某一次起以后全是反面”是同一件事。
\end{enumerate}
上极限的定义 $\bigcap_n\bigcup_{m\geqslant n}A_m$ 可以这样读：“无论 $n$ 多大，在第 $n$ 次以后总还会有 $A_m$ 发生”。这和数列的上极限是完全平行的，命题中的 (4) 正说明了这一点。以后的 Borel--Cantelli 引理就是在研究 $\P(\limsup_nA_n)$ 什么时候等于 $0$、什么时候等于 $1$。
\end{remark}

\begin{workedexample}\label{we:set-systems:limsup}
设 $n$ 为奇数时 $A_n=[0,1-\frac1n]$，$n$ 为偶数时 $A_n=[0,1+\frac1n]$。求 $\limsup_nA_n$ 与 $\liminf_nA_n$。
\end{workedexample}
\begin{solution}
我们用命题~\ref{prop:set-systems:limsup} 的 (1)(2)，对每个点 $x$ 判断它属于多少个 $A_n$。
\begin{itemize}
\item $x<0$ 或 $x>\frac32$：不属于任何 $A_n$（$A_n\subseteq[0,\frac32]$）。
\item $0\leqslant x<1$：取 $N$ 使 $\frac1N<1-x$，则 $n\geqslant N$ 时 $x\leqslant1-\frac1n$，从而 $x$ 属于一切 $A_n$（$n\geqslant N$），所以 $x\in\liminf_nA_n$。
\item $x=1$：$1$ 属于每个偶数项 $A_n$，不属于任何奇数项。属于无穷多个，但不是从某项起全属于。所以 $1\in\limsup_nA_n$，$1\notin\liminf_nA_n$。
\item $1<x\leqslant\frac32$：$x\in A_n$ 要求 $n$ 为偶数且 $\frac1n\geqslant x-1$，即 $n\leqslant\frac1{x-1}$，只有有限个 $n$。
\end{itemize}
所以 $\limsup_nA_n=[0,1]$，$\liminf_nA_n=[0,1)$。两者只差一个点 $1$，而这个点恰好是“奇偶两种情形来回摆动”造成的。
\end{solution}

对递增的集合列，记 $A_n\uparrow A$，表示 $A_1\subseteq A_2\subseteq\cdots$ 且 $\bigcup_nA_n=A$；对递减的集合列，记 $A_n\downarrow A$，表示 $A_1\supseteq A_2\supseteq\cdots$ 且 $\bigcap_nA_n=A$。

\topic{代数与 \texorpdfstring{$\sigma$}{σ}-代数}

\begin{definition}[代数]\label{def:set-systems:algebra}
设 $\cA\subseteq2^\Omega$。若
\begin{enumerate}
\item $\Omega\in\cA$；
\item 若 $A\in\cA$，则 $A^c\in\cA$；
\item 若 $A,B\in\cA$，则 $A\cup B\in\cA$，
\end{enumerate}
则称 $\cA$ 为 $\Omega$ 上的一个\term{代数}（algebra），也称\term{域}（field）。
\end{definition}

\begin{definition}[$\sigma$-代数]\label{def:set-systems:sigma-algebra}
设 $\cF\subseteq2^\Omega$。若
\begin{enumerate}
\item $\Omega\in\cF$；
\item 若 $A\in\cF$，则 $A^c\in\cF$；
\item 若 $A_1,A_2,\dots\in\cF$，则 $\bigcup_{n=1}^\infty A_n\in\cF$，
\end{enumerate}
则称 $\cF$ 为 $\Omega$ 上的一个 \term{$\sigma$-代数}（$\sigma$-algebra），也称 \term{$\sigma$-域}（$\sigma$-field）。二元组 $(\Omega,\cF)$ 称为\term{可测空间}（measurable space），$\cF$ 中的元素称为\term{可测集}（measurable set），在概率论中称为\term{事件}（event）。
\end{definition}

有些书的定义只写 (2)、(3) 两条，另外要求 $\cF$ 非空。二者等价：若 $\cF$ 非空，任取 $A\in\cF$，则 $A^c\in\cF$，于是 $\Omega=A\cup A^c\cup A\cup\cdots\in\cF$。

\begin{remark}[为什么是“可数”]
代数与 $\sigma$-代数的区别只在于允许有限并还是可数并。我们指出，“可数”恰好是一个合适的折中：
\begin{enumerate}
\item 只允许有限并是不够的。取极限是分析的基本操作，而像 $\limsup_nA_n$ 这样的集合一定要用到可数次运算。代数上当然也可以讨论“有限可加”的概率，但下一讲会看到，这样得到的“概率”连最基本的极限都无法交换。
\item 允许任意并又太多了。在 $[0,1]$ 上随机取一个数，每个单点的概率都是 $0$，而 $[0,1]$ 是不可数个单点的并，概率却是 $1$。如果要求概率对任意多个不交集合可加，这就会导出矛盾。
\end{enumerate}
顺便说一句，$\sigma$-代数中的 $\sigma$ 来自德文 Summe 的首字母，当时人们把并集称为和，$\sigma$ 表示的就是可数并。
\end{remark}

\begin{proposition}[代数与 $\sigma$-代数的封闭性]\label{prop:set-systems:closure}
\begin{enumerate}
\item 设 $\cA$ 是代数。则 $\varnothing\in\cA$，并且 $\cA$ 对有限交、有限并、差集与对称差封闭。
\item 设 $\cF$ 是 $\sigma$-代数。则 $\cF$ 是代数，并且对可数交封闭；若 $A_n\in\cF$，则 $\limsup_nA_n\in\cF$，$\liminf_nA_n\in\cF$。
\end{enumerate}
\end{proposition}
\begin{proof}
(1) $\varnothing=\Omega^c\in\cA$。若 $A,B\in\cA$，则 $A\cap B=(A^c\cup B^c)^c\in\cA$，这里依次用了补集封闭、并集封闭、补集封闭。对 $n$ 归纳，$A_1\cup\cdots\cup A_n=(A_1\cup\cdots\cup A_{n-1})\cup A_n\in\cA$，有限交同理。差集 $A\setminus B=A\cap B^c\in\cA$，对称差 $A\triangle B=(A\setminus B)\cup(B\setminus A)\in\cA$。

(2) 若 $A,B\in\cF$，令 $A_1=A$，$A_n=B$（$n\geqslant2$），则 $A\cup B=\bigcup_nA_n\in\cF$，故 $\cF$ 是代数。若 $A_n\in\cF$，则 $\bigcap_nA_n=(\bigcup_nA_n^c)^c\in\cF$。对每个 $n$，$\bigcup_{m\geqslant n}A_m\in\cF$，再对 $n$ 取可数交得 $\limsup_nA_n\in\cF$；$\liminf_nA_n$ 同理。
\end{proof}

\begin{example}[最简单的 $\sigma$-代数]
$\{\varnothing,\Omega\}$ 是 $\Omega$ 上最小的 $\sigma$-代数，称为\term{平凡 $\sigma$-代数}（trivial $\sigma$-algebra）；$2^\Omega$ 是最大的 $\sigma$-代数。对任意 $A\subseteq\Omega$，$\{\varnothing,A,A^c,\Omega\}$ 是 $\sigma$-代数：它只有有限个元素，对补集封闭是显然的，任意一列元素的并只能是这四个集合之一。
\end{example}

\begin{example}[有限划分生成的 $\sigma$-代数]\label{ex:set-systems:partition}
设 $\Omega=B_1+B_2+\cdots+B_k$ 是 $\Omega$ 的一个划分，$B_i$ 非空且两两不交。令
\[
\cF=\Bigl\{\bigcup_{i\in I}B_i:\ I\subseteq\{1,\dots,k\}\Bigr\}.
\]
则 $\cF$ 是 $\sigma$-代数，恰有 $2^k$ 个元素。事实上，$I=\{1,\dots,k\}$ 给出 $\Omega$；$\bigcup_{i\in I}B_i$ 的补集是 $\bigcup_{i\notin I}B_i$，这是因为各块两两不交且并为 $\Omega$；一列形如 $\bigcup_{i\in I_n}B_i$ 的集合的并是 $\bigcup_{i\in\bigcup_nI_n}B_i$。不同的 $I$ 给出不同的集合，因为各块非空且不交，所以共有 $2^k$ 个元素。可数划分 $\Omega=\sum_{i=1}^\infty B_i$ 的情形完全类似，此时 $I$ 取遍 $\N$ 的一切子集。

这个例子有很直观的概率意义。设想我们只能观察到结果落在哪一块，而看不到结果本身，那么我们能够判断发生与否的事件，恰好是 $\cF$ 中的集合。比如说，把一枚硬币抛两次，$\Omega=\{\text{正正},\text{正反},\text{反正},\text{反反}\}$。若只知道第一次的结果，划分就是 $\{\text{正正},\text{正反}\}$ 与 $\{\text{反正},\text{反反}\}$，对应的 $\sigma$-代数有 $4$ 个元素，“第一次是正面”在其中，“两次结果相同”不在其中。划分越细，$\sigma$-代数越大，我们掌握的信息就越多。以后讨论条件期望和鞅的时候，我们总是把 $\sigma$-代数理解为信息，这个例子是理解这种说法的最好的出发点。
\end{example}

\begin{example}[代数但不是 $\sigma$-代数]\label{ex:set-systems:finite-cofinite}
设 $\Omega=\N$，令 $\cA=\{A\subseteq\N: A\ \text{有限，或}\ A^c\ \text{有限}\}$。则 $\cA$ 是代数：$\N^c=\varnothing$ 有限，故 $\N\in\cA$；对补集封闭由定义的对称性立得；若 $A,B\in\cA$，当 $A,B$ 都有限时 $A\cup B$ 有限，否则不妨设 $A^c$ 有限，则 $(A\cup B)^c\subseteq A^c$ 有限。但 $\cA$ 不是 $\sigma$-代数：单点集 $\{2n\}$ 都属于 $\cA$，它们的并是偶数集，偶数集和它的补集（奇数集）都是无限集。
\end{example}

\begin{example}[可数－余可数 $\sigma$-代数]\label{ex:set-systems:countable-cocountable}
设 $\Omega$ 是任意集合，令 $\cF=\{A\subseteq\Omega: A\ \text{可数，或}\ A^c\ \text{可数}\}$（这里“可数”包括有限）。则 $\cF$ 是 $\sigma$-代数。$\Omega\in\cF$ 与补集封闭同上一例。设 $A_n\in\cF$。若每个 $A_n$ 都可数，则 $\bigcup_nA_n$ 是可数个可数集的并，仍可数；若某个 $A_{n_0}^c$ 可数，则 $(\bigcup_nA_n)^c\subseteq A_{n_0}^c$ 可数。当 $\Omega$ 不可数时，$\cF\neq2^\Omega$；比如 $\Omega=\R$ 时，$[0,\infty)$ 和它的补集都不可数，故 $[0,\infty)\notin\cF$。
\end{example}

\topic{生成的 \texorpdfstring{$\sigma$}{σ}-代数}

\begin{proposition}\label{prop:set-systems:intersection}
设 $I$ 是任意非空指标集（可以不可数），$\{\cF_i\}_{i\in I}$ 是 $\Omega$ 上的一族 $\sigma$-代数。则 $\bigcap_{i\in I}\cF_i=\{A\subseteq\Omega: A\in\cF_i\ \text{对一切}\ i\in I\ \text{成立}\}$ 是 $\sigma$-代数。对代数也有同样的结论。
\end{proposition}
\begin{proof}
记 $\cF=\bigcap_i\cF_i$。每个 $\cF_i$ 都含 $\Omega$，故 $\Omega\in\cF$。若 $A\in\cF$，则对每个 $i$，$A\in\cF_i$，从而 $A^c\in\cF_i$，故 $A^c\in\cF$。若 $A_n\in\cF$，则对每个 $i$，所有 $A_n$ 都在 $\cF_i$ 中，从而 $\bigcup_nA_n\in\cF_i$，故 $\bigcup_nA_n\in\cF$。代数的情形只需把可数并换成两个集合的并。
\end{proof}

注意两个 $\sigma$-代数的并一般不是 $\sigma$-代数。例如 $\Omega=\{1,2,3\}$，$\cF_1=\{\varnothing,\{1\},\{2,3\},\Omega\}$，$\cF_2=\{\varnothing,\{2\},\{1,3\},\Omega\}$，则 $\{1\},\{2\}\in\cF_1\cup\cF_2$，但 $\{1,2\}\notin\cF_1\cup\cF_2$。

\begin{definition}[生成的 $\sigma$-代数]\label{def:set-systems:generated}
设 $\cC\subseteq2^\Omega$ 是任意一个集合族。所有包含 $\cC$ 的 $\sigma$-代数的交
\[
\sigma(\cC)=\bigcap\{\cF:\ \cF\ \text{是}\ \Omega\ \text{上的}\ \sigma\text{-代数},\ \cC\subseteq\cF\}
\]
称为由 $\cC$ \term{生成的 $\sigma$-代数}（$\sigma$-algebra generated by $\cC$）。
\end{definition}

由于 $2^\Omega$ 是包含 $\cC$ 的 $\sigma$-代数，上面的交是对一个非空族取的，由命题~\ref{prop:set-systems:intersection}，$\sigma(\cC)$ 是 $\sigma$-代数。它包含 $\cC$，并且含于任何包含 $\cC$ 的 $\sigma$-代数之中，所以 $\sigma(\cC)$ 是包含 $\cC$ 的\uline{最小}的 $\sigma$-代数。同样可以定义由 $\cC$ 生成的代数，记作 $\alpha(\cC)$。

\begin{proposition}\label{prop:set-systems:generated}
设 $\cC,\cD\subseteq2^\Omega$。
\begin{enumerate}
\item 若 $\cC$ 本身是 $\sigma$-代数，则 $\sigma(\cC)=\cC$。
\item 若 $\cC\subseteq\cD$，则 $\sigma(\cC)\subseteq\sigma(\cD)$。
\item 若 $\cC\subseteq\sigma(\cD)$，则 $\sigma(\cC)\subseteq\sigma(\cD)$。特别地，若 $\cC\subseteq\sigma(\cD)$ 且 $\cD\subseteq\sigma(\cC)$，则 $\sigma(\cC)=\sigma(\cD)$。
\end{enumerate}
\end{proposition}
\begin{proof}
(1) $\cC$ 是包含 $\cC$ 的 $\sigma$-代数，所以 $\sigma(\cC)\subseteq\cC$；反包含由定义给出。(2) $\sigma(\cD)$ 是包含 $\cD$ 从而包含 $\cC$ 的 $\sigma$-代数，由最小性 $\sigma(\cC)\subseteq\sigma(\cD)$。(3) $\sigma(\cD)$ 是包含 $\cC$ 的 $\sigma$-代数，由最小性即得第一句；第二句对两个方向各用一次第一句。
\end{proof}

\begin{remark}
生成的 $\sigma$-代数是“从外面”定义的：我们取了一切包含 $\cC$ 的 $\sigma$-代数的交，而并没有写出 $\sigma(\cC)$ 里的元素长什么样。请比较线性代数中由一组向量张成的子空间：那里既可以说“包含这组向量的最小子空间”，也可以直接写出一般元素，即有限线性组合。对 $\sigma$-代数，后一种描述很难做到。如果试图从 $\cC$ 出发，反复地取补集、取可数并，把得到的集合都放进来，那么做一次、两次、可数多次都不够，要一直做“不可数多次”才能停下来（本讲作业 F）。所以对 $\sigma(\cC)$，我们几乎从不去考察一般元素的具体形状。
\end{remark}

这种定义方式既是方便之处，也是困难之处。要证明“$\sigma(\cC)$ 中的每个集合都具有性质 $(\ast)$”，一般的办法是下面的\term{好集原理}（good sets principle）：
\begin{quote}
令 $\cG=\{A\in\sigma(\cC): A\ \text{具有性质}\ (\ast)\}$（“好集”的全体）。若能证明 (a) $\cC\subseteq\cG$，(b) $\cG$ 是 $\sigma$-代数，则由最小性 $\sigma(\cC)\subseteq\cG$，即 $\sigma(\cC)$ 中每个集合都是好集。
\end{quote}
这个原理用起来的困难在于步骤 (b)：很多自然的性质，我们只能验证它对补集和递增并封闭，却无法验证它对交集封闭。本讲最后一节的 $\pi$-$\lambda$ 定理就是用来解决这个困难的。

\begin{example}[单点集生成的 $\sigma$-代数]
设 $\Omega$ 是任意集合，$\cC=\{\{\omega\}:\omega\in\Omega\}$。则 $\sigma(\cC)$ 就是\egref{ex:set-systems:countable-cocountable} 中的可数－余可数 $\sigma$-代数 $\cF$。一方面，$\cF$ 是包含全体单点集的 $\sigma$-代数，所以 $\sigma(\cC)\subseteq\cF$。另一方面，若 $A$ 可数，则 $A$ 是可数个单点集的并，从而 $A\in\sigma(\cC)$；若 $A^c$ 可数，则 $A^c\in\sigma(\cC)$，从而 $A\in\sigma(\cC)$。于是 $\cF\subseteq\sigma(\cC)$。

这个例子说明，生成元选得不好，得到的 $\sigma$-代数可能非常小：在 $\R$ 上，仅仅知道每个单点是事件，得到的 $\sigma$-代数里连 $[0,\infty)$ 都没有。
\end{example}

\begin{workedexample}\label{we:set-systems:two-sets}
设 $A,B\subseteq\Omega$。写出 $\sigma(\{A,B\})$ 的全部元素。
\end{workedexample}
\begin{solution}
思路：先找出 $A,B$ 把 $\Omega$ 切成的最小的块，再用\egref{ex:set-systems:partition}。

令
\[
B_1=A\cap B,\quad B_2=A\cap B^c,\quad B_3=A^c\cap B,\quad B_4=A^c\cap B^c .
\]
每个 $\omega$ 恰好属于其中一个（看 $\omega$ 是否属于 $A$、是否属于 $B$，共四种情形），所以它们两两不交，并为 $\Omega$。去掉其中的空集，剩下的构成 $\Omega$ 的一个划分。令 $\cF$ 为若干块之并的全体。

\emph{$\cF\subseteq\sigma(\{A,B\})$。}每个 $B_i$ 是 $A$ 或 $A^c$ 与 $B$ 或 $B^c$ 的交，属于 $\sigma(\{A,B\})$（命题~\ref{prop:set-systems:closure}），它们的并也属于。

\emph{$\sigma(\{A,B\})\subseteq\cF$。}由\egref{ex:set-systems:partition}，$\cF$ 是 $\sigma$-代数；又 $A=B_1\cup B_2\in\cF$，$B=B_1\cup B_3\in\cF$。由最小性得 $\sigma(\{A,B\})\subseteq\cF$。

所以 $\sigma(\{A,B\})=\cF$，当四块都非空时它恰有 $2^4=16$ 个元素，比如 $A\cup B=B_1\cup B_2\cup B_3$，$A\triangle B=B_2\cup B_3$。若 $A\subseteq B$，则 $B_2=\varnothing$，只剩三块，$\sigma(\{A,B\})$ 有 $8$ 个元素。

同样的方法说明：$n$ 个集合 $A_1,\dots,A_n$ 把 $\Omega$ 切成至多 $2^n$ 块，形如 $A_1^{\pm}\cap\cdots\cap A_n^{\pm}$（$A^+=A$，$A^-=A^c$），$\sigma(\{A_1,\dots,A_n\})$ 就是这些块的并的全体，至多有 $2^{2^n}$ 个元素。有限个集合生成的 $\sigma$-代数总是有限的，并且可以完全写出来；到了可数个生成元，情形就完全不同了，见下面的例子和本讲作业 C。
\end{solution}

\begin{example}[无穷次抛硬币的 $\sigma$-代数]\label{ex:set-systems:coin-tossing}
回到本讲开头的 $\Omega=\{0,1\}^{\N}$。令
\[
A_n=\{\omega\in\Omega:\ \omega_n=1\}
\]
为“第 $n$ 次是正面”。由前 $n$ 次抛掷决定的事件组成的 $\sigma$-代数是 $\cF_n=\sigma(\{A_1,\dots,A_n\})$。按上一个例题，它由 $2^n$ 个块
\[
\{\omega:\ \omega_1=\varepsilon_1,\dots,\omega_n=\varepsilon_n\},\qquad \varepsilon_1,\dots,\varepsilon_n\in\{0,1\}
\]
的并组成（这些块叫做\term{柱集}（cylinder set），因为只限制了前 $n$ 个坐标）。比如“前三次中恰有两次正面”属于 $\cF_3$，是 $3$ 个柱集的并。$\cF_1\subseteq\cF_2\subseteq\cdots$，这表示随着抛掷的进行，我们掌握的信息越来越多。

整个试验的 $\sigma$-代数取为 $\cF=\sigma(\{A_n:n\in\N\})$。每个 $\cF_n$ 都含于 $\cF$，但 $\cF$ 比 $\bigcup_n\cF_n$ 大得多：$\bigcup_n\cF_n$ 中的每个事件只依赖于有限次抛掷，而“正面出现无穷多次”$=\limsup_nA_n\in\cF$ 依赖于全部的抛掷，它不属于任何一个 $\cF_n$（因为改变前 $n$ 个坐标不会改变它是否发生，而 $\cF_n$ 中除 $\varnothing,\Omega$ 外的集合都会被这种改变影响）。本讲作业习题~\exref{ex:set-systems:A4} 中有类似的现象。
\end{example}

\begin{workedexample}\label{we:set-systems:slln-event}
在上例中，证明 $L=\{\omega:\ \lim_{n\to\infty}S_n(\omega)/n=\frac12\}\in\cF$，其中 $S_n=\omega_1+\cdots+\omega_n$。
\end{workedexample}
\begin{solution}
思路：把“数列收敛到 $\frac12$”的 $\varepsilon$-$N$ 语言逐字翻译成集合运算。“对一切”翻译成交，“存在”翻译成并；为了只用到可数次运算，把 $\varepsilon$ 限制为 $\frac1k$。

\emph{第一步：翻译。}$S_n(\omega)/n\to\frac12$ 的意思是：对每个 $k\in\N$，存在 $N\in\N$，使一切 $n\geqslant N$ 有 $\abs{S_n(\omega)/n-\frac12}\leqslant\frac1k$。（用 $\varepsilon=\frac1k$ 代替一切 $\varepsilon>0$ 是可以的：对任意 $\varepsilon>0$ 取 $k>\frac1\varepsilon$ 即可。）令
\[
E_{n,k}=\Bigl\{\omega:\ \Bigl\lvert\frac{S_n(\omega)}n-\frac12\Bigr\rvert\leqslant\frac1k\Bigr\},
\]
则
\[
L=\bigcap_{k=1}^\infty\ \bigcup_{N=1}^\infty\ \bigcap_{n=N}^\infty E_{n,k}.
\]
等式的两边逐字对应：$\omega\in L$ 当且仅当对每个 $k$（交），存在 $N$（并），使对每个 $n\geqslant N$（交）都有 $\omega\in E_{n,k}$。

\emph{第二步：$E_{n,k}\in\cF$。}$S_n(\omega)$ 只依赖于前 $n$ 个坐标，所以 $E_{n,k}$ 是若干个柱集 $\{\omega_1=\varepsilon_1,\dots,\omega_n=\varepsilon_n\}$ 的并，并的范围是满足 $\abs{(\varepsilon_1+\cdots+\varepsilon_n)/n-\frac12}\leqslant\frac1k$ 的那些 $(\varepsilon_1,\dots,\varepsilon_n)$。每个柱集等于 $A_1^{\pm}\cap\cdots\cap A_n^{\pm}$，属于 $\cF_n\subseteq\cF$。所以 $E_{n,k}\in\cF_n\subseteq\cF$。

\emph{第三步。}$\cF$ 对可数交、可数并封闭（命题~\ref{prop:set-systems:closure}），由第一步的表达式，$L\in\cF$。
\end{solution}

这个例题的做法以后会反复出现：一个关于极限的事件，按照极限的定义写成“交、并、交”的形式，每个最内层的集合只涉及有限个坐标，从而是事件。它也说明了，如果 $\sigma$-代数只对有限次运算封闭，我们连强大数定律的陈述都写不出来。

\begin{remark}[$\sigma$-代数与信息]
在\egref{ex:set-systems:coin-tossing} 中，$\cF_n$ 由前 $n$ 次抛掷决定。一个事件 $A$ 属于 $\cF_n$，直观上就是说“看了前 $n$ 次的结果，就能判断 $A$ 发生了没有”。比如 $A_1\cap A_2$ 属于 $\cF_2$；“第 $3$ 次是正面”不属于 $\cF_2$；$\limsup_nA_n$ 不属于任何 $\cF_n$，看了任何有限次的结果都无法判断。这种“$\sigma$-代数就是信息”的观点在本课程后半部分讲鞅的时候是基本的。
\end{remark}

\topic{Borel \texorpdfstring{$\sigma$}{σ}-代数}

\begin{definition}[Borel $\sigma$-代数]\label{def:set-systems:borel}
设 $S$ 是拓扑空间（例如度量空间），由 $S$ 的全体开集生成的 $\sigma$-代数称为 $S$ 上的 \term{Borel $\sigma$-代数}（Borel $\sigma$-algebra），记作 $\cB(S)$，其中的元素称为 \term{Borel 集}（Borel set）。$\R^d$ 上的 Borel $\sigma$-代数记作 $\cR^d$，$d=1$ 时简记为 $\cR$。
\end{definition}

由于闭集是开集的补集，$\cB(S)$ 也由全体闭集生成。$\cR$ 中包含了一切开集、闭集、可数集（单点集 $\{x\}$ 是闭集）、区间，以及对它们作可数次运算得到的集合，例如全体无理数 $\R\setminus\Q$。下面的命题说明 $\cR$ 可以由很小的集合族生成。

\begin{lemma}\label{lem:set-systems:open-sets-R}
$\R$ 中的每个开集都是可数个端点为有理数的开区间的并。
\end{lemma}
\begin{proof}
设 $U\subseteq\R$ 是开集。令 $\cI=\{(p,q): p,q\in\Q,\ p<q,\ (p,q)\subseteq U\}$。$\cI$ 是可数集 $\Q\times\Q$ 的子集的像，故可数，并且 $\bigcup\cI\subseteq U$。反之，任取 $x\in U$，由于 $U$ 是开集，存在 $\varepsilon>0$ 使 $(x-\varepsilon,x+\varepsilon)\subseteq U$。由有理数的稠密性，取有理数 $p\in(x-\varepsilon,x)$，$q\in(x,x+\varepsilon)$，则 $x\in(p,q)\subseteq(x-\varepsilon,x+\varepsilon)\subseteq U$，即 $(p,q)\in\cI$。所以 $U=\bigcup\cI$。
\end{proof}

\begin{proposition}\label{prop:set-systems:borel-generators}
下列每个集合族都生成 $\cR$：
\begin{enumerate}
\item $\cC_1=\{(a,b): a<b,\ a,b\in\R\}$；
\item $\cC_2=\{(a,b]: a<b,\ a,b\in\R\}$；
\item $\cC_3=\{(-\infty,x]: x\in\R\}$；
\item $\cC_4=\{(-\infty,x]: x\in\Q\}$。
\end{enumerate}
\end{proposition}
\begin{proof}
我们证明 $\cR=\sigma(\cC_1)=\sigma(\cC_2)=\sigma(\cC_3)=\sigma(\cC_4)$，每一步都用命题~\ref{prop:set-systems:generated}(3)：只需验证每个集合族包含于另一个生成的 $\sigma$-代数中。

\emph{$\sigma(\cC_1)=\cR$。}开区间是开集，故 $\cC_1\subseteq\cR$，从而 $\sigma(\cC_1)\subseteq\cR$。由引理~\ref{lem:set-systems:open-sets-R}，每个开集都是 $\cC_1$ 中可数个集合的并，属于 $\sigma(\cC_1)$，所以 $\cR\subseteq\sigma(\cC_1)$。

\emph{$\sigma(\cC_2)=\sigma(\cC_1)$。}由
\[
(a,b]=\bigcap_{n=1}^\infty\Bigl(a,b+\frac1n\Bigr),\qquad (a,b)=\bigcup_{n>1/(b-a)}\Bigl(a,b-\frac1n\Bigr]
\]
知 $\cC_2\subseteq\sigma(\cC_1)$，$\cC_1\subseteq\sigma(\cC_2)$。第一个等式成立是因为 $x\in(a,b+1/n)$ 对一切 $n$ 成立当且仅当 $a<x\leqslant b$；第二个等式中 $n>1/(b-a)$ 保证了 $b-1/n>a$，而 $a<x<b$ 当且仅当对某个充分大的 $n$ 有 $a<x\leqslant b-1/n$。

\emph{$\sigma(\cC_3)=\sigma(\cC_2)$。}由 $(a,b]=(-\infty,b]\setminus(-\infty,a]$ 知 $\cC_2\subseteq\sigma(\cC_3)$；由 $(-\infty,x]=\bigcup_{n=1}^\infty(x-n,x]$ 知 $\cC_3\subseteq\sigma(\cC_2)$。

\emph{$\sigma(\cC_4)=\sigma(\cC_3)$。}$\cC_4\subseteq\cC_3$ 给出 $\sigma(\cC_4)\subseteq\sigma(\cC_3)$。对任意 $x\in\R$，取有理数列 $q_n\downarrow x$ 且 $q_n>x$（例如在 $(x,x+1/n)$ 中取有理数，再令 $q_n$ 为前 $n$ 个所取之数的最小值），则 $(-\infty,x]=\bigcap_n(-\infty,q_n]$：若 $y\leqslant x$，则 $y\leqslant q_n$ 对一切 $n$ 成立；若 $y>x$，由于 $q_n\to x$，存在 $n$ 使 $q_n<y$。故 $\cC_3\subseteq\sigma(\cC_4)$。
\end{proof}

\begin{remark}
命题说明，要验证某个集合族包含 $\cR$，只要验证它包含一切 $(-\infty,x]$ 就可以了。这是今后反复使用的原则：随机变量 $X$ 的分布由 $\P(X\leqslant x)$ 决定，$X$ 是否可测只要看 $\{X\leqslant x\}$ 是不是事件，都源于此。$\cR$ 中的集合非常多，我们在分析和概率中能够“写出来”的集合，只要不借助选择公理，基本上都是 Borel 集；但 $\cR$ 并不是 $\R$ 的全体子集（本讲作业 F，以及第~\ref{lec:extension-of-measures} 讲的作业）。
\end{remark}

\begin{example}[几个具体的 Borel 集]
\begin{enumerate}
\item $\Q=\bigcup_{q\in\Q}\{q\}$ 是可数个闭集的并，是 Borel 集；无理数集 $\R\setminus\Q$ 是它的补集，也是 Borel 集。
\item \emph{Cantor 集。}令 $C_0=[0,1]$，$C_1=[0,\frac13]\cup[\frac23,1]$，一般地 $C_{n+1}$ 是把 $C_n$ 的每个闭区间三等分后去掉中间开区间所得。每个 $C_n$ 是有限个闭区间的并，是闭集。Cantor 集 $C=\bigcap_nC_n$ 是可数个闭集的交，是 Borel 集。
\item \emph{二进制数字。}对 $x\in[0,1)$，写 $x=\sum_{n\geqslant1}d_n(x)2^{-n}$，约定不取以全 $1$ 结尾的展开。$\{x: d_n(x)=1\}$ 是 $2^{n-1}$ 个形如 $[\frac{2j+1}{2^n},\frac{2j+2}{2^n})$ 的区间的并，是 Borel 集。仿照例题（第~\pageref{we:set-systems:slln-event} 页）的写法，$\{x:\ \frac1n\sum_{i\leqslant n}d_i(x)\to\frac12\}$ 也是 Borel 集。第~\ref{lec:independence} 讲的 Borel 正规数定理要说明，这个集合的 Lebesgue 测度等于 $1$。
\end{enumerate}
\end{example}

\begin{workedexample}\label{we:set-systems:translation}
设 $B\in\cR$，$c\in\R$，$a\neq0$。证明 $B+c=\{x+c:x\in B\}$ 与 $aB=\{ax:x\in B\}$ 都是 Borel 集。
\end{workedexample}
\begin{solution}
思路：我们无法写出一般 Borel 集的样子，所以不能直接验证。用好集原理：把具有所要性质的集合收集起来，证明它们构成 $\sigma$-代数并且包含开集。

令 $\cG=\{B\in\cR:\ B+c\in\cR\}$。
\begin{itemize}
\item \emph{$\cG$ 包含一切开集。}若 $U$ 是开集，$y=x+c\in U+c$，$x\in U$，取 $\varepsilon$ 使 $(x-\varepsilon,x+\varepsilon)\subseteq U$，则 $(y-\varepsilon,y+\varepsilon)\subseteq U+c$。所以 $U+c$ 是开集，属于 $\cR$。
\item \emph{$\R\in\cG$}：$\R+c=\R$。
\item \emph{补集。}平移是双射，所以 $B^c+c=(B+c)^c$。若 $B\in\cG$，则 $(B+c)^c\in\cR$，即 $B^c\in\cG$。
\item \emph{可数并。}$(\bigcup_nB_n)+c=\bigcup_n(B_n+c)$。若 $B_n\in\cG$，右边是可数个 Borel 集的并，故 $\bigcup_nB_n\in\cG$。
\end{itemize}
所以 $\cG$ 是包含全体开集的 $\sigma$-代数，由 $\cR$ 的最小性，$\cR\subseteq\cG$，即一切 Borel 集平移后仍是 Borel 集。

对 $aB$ 完全一样：$x\mapsto ax$ 是双射，把开集映成开集（$(x-\varepsilon,x+\varepsilon)$ 映成以 $ax$ 为中心、半径 $\abs a\varepsilon$ 的开区间），并保持补集与可数并。注意 $a\neq0$ 是必要的：$a=0$ 时这个映射不是双射，$B^c$ 的像不再是 $B$ 的像的补集。
\end{solution}

这个例题的证明方法同学们应该熟记：要证明\emph{每个} Borel 集都有某种性质，就证明“有这种性质的集合”构成 $\sigma$-代数，并且包含开集（或者包含一个更小的生成元族）。

在 $\R^d$ 上可以作完全平行的讨论：$\cR^d$ 由一切形如 $(a_1,b_1)\times\cdots\times(a_d,b_d)$ 的开矩形生成，也由形如 $(a_1,b_1]\times\cdots\times(a_d,b_d]$ 的矩形生成，还由 $\{(-\infty,x_1]\times\cdots\times(-\infty,x_d]: x_i\in\Q\}$ 生成。证明留作作业（习题~\exref{ex:set-systems:A5}）。

\topic{\texorpdfstring{$\pi$}{π}-类与 \texorpdfstring{$\lambda$}{λ}-类}

\begin{definition}[$\pi$-类]\label{def:set-systems:pi}
设 $\cP\subseteq2^\Omega$。若 $A,B\in\cP$ 蕴含 $A\cap B\in\cP$，则称 $\cP$ 为 \term{$\pi$-类}（$\pi$-system）。
\end{definition}

\begin{definition}[$\lambda$-类]\label{def:set-systems:lambda}
设 $\cL\subseteq2^\Omega$。若
\begin{enumerate}
\item $\Omega\in\cL$；
\item 若 $A,B\in\cL$ 且 $A\subseteq B$，则 $B\setminus A\in\cL$；
\item 若 $A_n\in\cL$ 且 $A_n\uparrow A$，则 $A\in\cL$，
\end{enumerate}
则称 $\cL$ 为 \term{$\lambda$-类}（$\lambda$-system），也称 \term{Dynkin 类}（Dynkin system）。
\end{definition}

\begin{example}[$\pi$-类]\label{ex:set-systems:pi-examples}
\begin{enumerate}
\item $\{(-\infty,x]: x\in\R\}$ 是 $\pi$-类，因为 $(-\infty,x]\cap(-\infty,y]=(-\infty,\min\{x,y\}]$。
\item $\{(a,b]: a<b\}\cup\{\varnothing\}$ 是 $\pi$-类：两个左开右闭区间的交是 $(\max\{a,a'\},\min\{b,b'\}]$，当 $\max\{a,a'\}\geqslant\min\{b,b'\}$ 时为空集。同理，$\R^d$ 中形如 $(a_1,b_1]\times\cdots\times(a_d,b_d]$ 的矩形连同空集组成 $\pi$-类。
\item 任何代数都是 $\pi$-类。
\end{enumerate}
\end{example}

\begin{proposition}[$\lambda$-类的等价定义]\label{prop:set-systems:lambda-equivalent}
$\cL\subseteq2^\Omega$ 是 $\lambda$-类，当且仅当
\begin{enumerate}
\item[(1$'$)] $\Omega\in\cL$；
\item[(2$'$)] 若 $A\in\cL$，则 $A^c\in\cL$；
\item[(3$'$)] 若 $A_1,A_2,\dots\in\cL$ 两两不交，则 $\bigcup_nA_n\in\cL$。
\end{enumerate}
\end{proposition}
\begin{proof}
设 $\cL$ 是 $\lambda$-类。(1$'$) 即 (1)。若 $A\in\cL$，由 $A\subseteq\Omega$ 及 (2) 得 $A^c=\Omega\setminus A\in\cL$，这是 (2$'$)。为证 (3$'$)，先证 $\cL$ 对两个不交集合的并封闭：设 $A,B\in\cL$，$A\cap B=\varnothing$，则 $A\subseteq B^c$，而 $B^c\in\cL$，由 (2) 得 $B^c\setminus A\in\cL$，再取补得 $(B^c\setminus A)^c=A\cup B\in\cL$。对 $n$ 归纳，若 $A_1,\dots,A_n\in\cL$ 两两不交，则 $B_n=A_1\cup\cdots\cup A_n\in\cL$（$B_{n-1}$ 与 $A_n$ 不交）。现在 $B_n\uparrow\bigcup_nA_n$，由 (3) 得 $\bigcup_nA_n\in\cL$。

反之，设 $\cL$ 满足 (1$'$)–(3$'$)。(1) 即 (1$'$)。设 $A,B\in\cL$，$A\subseteq B$。则 $B^c\in\cL$，且 $A$ 与 $B^c$ 不交，由 (3$'$)（取 $A_1=A$，$A_2=B^c$，其余为 $\varnothing=\Omega^c\in\cL$）得 $A\cup B^c\in\cL$，再取补得 $(A\cup B^c)^c=B\cap A^c=B\setminus A\in\cL$，这是 (2)。设 $A_n\in\cL$，$A_n\uparrow A$。令 $D_1=A_1$，$D_n=A_n\setminus A_{n-1}$（$n\geqslant2$），由刚证的 (2) 知 $D_n\in\cL$；它们两两不交（若 $m<n$，则 $D_m\subseteq A_m\subseteq A_{n-1}$，而 $D_n\cap A_{n-1}=\varnothing$），且 $\bigcup_{k\leqslant n}D_k=A_n$，从而 $\bigcup_nD_n=\bigcup_nA_n=A$。由 (3$'$) 得 $A\in\cL$，这是 (3)。
\end{proof}

\begin{example}[$\lambda$-类不一定是 $\sigma$-代数]\label{ex:set-systems:even}
设 $\Omega=\{1,2,3,4\}$，$\cL$ 是 $\Omega$ 的全体偶数元子集：
\[
\cL=\bigl\{\varnothing,\ \Omega,\ \{1,2\},\{3,4\},\{1,3\},\{2,4\},\{1,4\},\{2,3\}\bigr\}.
\]
用命题~\ref{prop:set-systems:lambda-equivalent} 验证：$\Omega$ 有 $4$ 个元素；偶数元子集的补集有 $4-\text{偶数}$ 个元素，仍为偶数；由于 $\Omega$ 有限，一列两两不交的非空集合只有有限个，而不交的偶数元集合的并仍是偶数元集合。所以 $\cL$ 是 $\lambda$-类。但 $\{1,2\}\cap\{1,3\}=\{1\}\notin\cL$，所以 $\cL$ 不是 $\pi$-类，更不是 $\sigma$-代数。
\end{example}

为什么要引入 $\lambda$-类这样一个看起来不太自然的概念？我们回到好集原理，看一个具体的问题。设 $\mu,\nu$ 是 $(\Omega,\cF)$ 上的两个概率（下一讲给出严格定义），我们想知道它们何时相等。令
\[
\cD=\{A\in\cF: \mu(A)=\nu(A)\}
\]
为“好集”的全体。利用概率的基本性质（下一讲的定理~\ref{thm:measures:basic}），我们可以检验：
\begin{enumerate}
\item $\Omega\in\cD$，因为 $\mu(\Omega)=\nu(\Omega)=1$；
\item 若 $A\subseteq B$ 都在 $\cD$ 中，则 $\mu(B\setminus A)=\mu(B)-\mu(A)=\nu(B)-\nu(A)=\nu(B\setminus A)$；
\item 若 $A_n\uparrow A$ 且 $A_n\in\cD$，则由概率的连续性，$\mu(A)=\lim\mu(A_n)=\lim\nu(A_n)=\nu(A)$。
\end{enumerate}
这恰好是 $\lambda$-类的三条公理。但是，从 $\mu(A)=\nu(A)$ 与 $\mu(B)=\nu(B)$ 推不出 $\mu(A\cap B)=\nu(A\cap B)$，所以我们没有办法说明 $\cD$ 是 $\sigma$-代数。比如说，在 $\Omega=\{1,2,3,4\}$ 上取 $\mu$ 为均匀分布（每点概率 $1/4$），$\nu$ 给 $1,4$ 各 $1/2$ 的概率，给 $2,3$ 概率 $0$。逐个比较 $16$ 个子集可知
\[
\cD=\bigl\{\varnothing,\ \Omega,\ \{1,2\},\{3,4\},\{1,3\},\{2,4\}\bigr\}.
\]
它是 $\lambda$-类，但 $\{1,2\}\cap\{1,3\}=\{1\}\notin\cD$。

\begin{remark}[$\pi$-类与 $\lambda$-类的分工]
上面的讨论不是偶然的。凡是用等式和极限定义出来的性质，比如两个测度相等、两个积分相等、某个函数等式对示性函数成立，好集的全体对差集和递增并封闭几乎是自动的，而对交集封闭通常没有理由成立。$\sigma$-代数的定义把两种要求捆在了一起，而 $\pi$-类与 $\lambda$-类把它们拆开了：
\begin{enumerate}
\item 对交集封闭的要求，交给生成元去验证，生成元是我们自己选的，可以选成 $\pi$-类，比如一切 $(-\infty,x]$；
\item 对差集和递增并封闭的要求，交给好集去验证，这通常由测度或积分的性质自动给出。
\end{enumerate}
下一节的 $\pi$-$\lambda$ 定理说，这两半合在一起就足够了。这就是这个定理在整个课程中反复出现的原因。
\end{remark}

\topic{Dynkin 的 \texorpdfstring{$\pi$-$\lambda$}{π-λ} 定理}

\begin{remark}[证明思路]
所要证明的是：包含 $\pi$-类 $\cP$ 的最小 $\lambda$-类 $\ell(\cP)$ 就是 $\sigma(\cP)$。证明分两大步：第一步，同时是 $\pi$-类和 $\lambda$-类的集合族一定是 $\sigma$-代数（引理~\ref{lem:set-systems:pi-lambda-sigma}）；第二步，$\ell(\cP)$ 仍是 $\pi$-类（引理~\ref{lem:set-systems:ell-is-pi}）。两步合起来，$\ell(\cP)$ 是包含 $\cP$ 的 $\sigma$-代数，从而包含 $\sigma(\cP)$。第一步是定义的直接验证，真正需要想法的是第二步。第一遍读时，可以先承认第二步，把定理~\ref{thm:set-systems:pi-lambda} 的陈述和后面的例题读懂。
\end{remark}

与命题~\ref{prop:set-systems:intersection} 的证明完全相同，任意一族 $\lambda$-类的交仍是 $\lambda$-类：$\Omega$ 属于每一个；若 $A\subseteq B$ 都属于每一个，则 $B\setminus A$ 属于每一个；若递增列 $A_n$ 的每一项都属于每一个，则 $\bigcup_nA_n$ 属于每一个。由于 $2^\Omega$ 是 $\lambda$-类，对任意 $\cC\subseteq2^\Omega$，可以定义包含 $\cC$ 的最小 $\lambda$-类
\[
\ell(\cC)=\bigcap\{\cL:\ \cL\ \text{是}\ \lambda\text{-类},\ \cC\subseteq\cL\}.
\]
由于 $\sigma$-代数一定是 $\lambda$-类（$B\setminus A=B\cap A^c$，递增并是可数并），$\sigma(\cC)$ 是包含 $\cC$ 的 $\lambda$-类，所以总有 $\ell(\cC)\subseteq\sigma(\cC)$。

\begin{lemma}[第一步]\label{lem:set-systems:pi-lambda-sigma}
若 $\cG\subseteq2^\Omega$ 既是 $\pi$-类又是 $\lambda$-类，则 $\cG$ 是 $\sigma$-代数。
\end{lemma}
\begin{proof}
逐条验证定义~\ref{def:set-systems:sigma-algebra}。

(1) $\Omega\in\cG$，因为 $\cG$ 是 $\lambda$-类。

(2) 若 $A\in\cG$，由 $A\subseteq\Omega$ 及 $\lambda$-类的性质 (2) 得 $A^c=\Omega\setminus A\in\cG$。

(3) 先证 $\cG$ 对有限并封闭。设 $A,B\in\cG$，由 (2) 得 $A^c,B^c\in\cG$，由 $\pi$-类的性质得 $A^c\cap B^c\in\cG$，再由 (2) 得 $A\cup B=(A^c\cap B^c)^c\in\cG$。归纳可得 $\cG$ 对任意有限并封闭。现设 $A_1,A_2,\dots\in\cG$，令 $B_n=A_1\cup\cdots\cup A_n$，则 $B_n\in\cG$，且 $B_n\uparrow\bigcup_{k=1}^\infty A_k$。由 $\lambda$-类的性质 (3) 得 $\bigcup_kA_k\in\cG$。
\end{proof}

\begin{lemma}[第二步]\label{lem:set-systems:ell-is-pi}
若 $\cP$ 是 $\pi$-类，则 $\ell(\cP)$ 是 $\pi$-类。
\end{lemma}
\begin{proof}
对每个 $A\subseteq\Omega$，令
\[
\cG_A=\{B\subseteq\Omega: A\cap B\in\ell(\cP)\}.
\]
我们要证明：对任意 $A,B\in\ell(\cP)$，都有 $B\in\cG_A$。分三步。

\emph{(a) 若 $A\in\ell(\cP)$，则 $\cG_A$ 是 $\lambda$-类。}逐条验证定义~\ref{def:set-systems:lambda}：
\begin{itemize}
\item $A\cap\Omega=A\in\ell(\cP)$，故 $\Omega\in\cG_A$。
\item 设 $B,C\in\cG_A$ 且 $C\subseteq B$。则 $A\cap C\subseteq A\cap B$，二者都属于 $\ell(\cP)$，由于 $\ell(\cP)$ 是 $\lambda$-类，$(A\cap B)\setminus(A\cap C)\in\ell(\cP)$。而
\[
(A\cap B)\setminus(A\cap C)=A\cap B\cap(A^c\cup C^c)=A\cap B\cap C^c=A\cap(B\setminus C),
\]
故 $B\setminus C\in\cG_A$。
\item 设 $B_n\in\cG_A$，$B_n\uparrow B$。则 $A\cap B_n\in\ell(\cP)$，且 $A\cap B_n\uparrow A\cap B$（递增是显然的，并且 $\bigcup_n(A\cap B_n)=A\cap\bigcup_nB_n$）。由 $\ell(\cP)$ 是 $\lambda$-类得 $A\cap B\in\ell(\cP)$，即 $B\in\cG_A$。
\end{itemize}

\emph{(b) 若 $A\in\cP$，则 $\ell(\cP)\subseteq\cG_A$。}对任意 $B\in\cP$，由于 $\cP$ 是 $\pi$-类，$A\cap B\in\cP\subseteq\ell(\cP)$，即 $B\in\cG_A$。所以 $\cP\subseteq\cG_A$。又 $A\in\cP\subseteq\ell(\cP)$，由 (a)，$\cG_A$ 是 $\lambda$-类。由 $\ell(\cP)$ 的最小性，$\ell(\cP)\subseteq\cG_A$。

\emph{(c) 若 $A\in\ell(\cP)$，则 $\ell(\cP)\subseteq\cG_A$。}任取 $B\in\cP$。由 (b)（把 (b) 中的 $A$ 换成 $B$），$\ell(\cP)\subseteq\cG_B$，而 $A\in\ell(\cP)$，所以 $A\in\cG_B$，即 $B\cap A\in\ell(\cP)$，也就是 $B\in\cG_A$。这说明 $\cP\subseteq\cG_A$。由 (a)，$\cG_A$ 是 $\lambda$-类，再由最小性得 $\ell(\cP)\subseteq\cG_A$。

由 (c)，对任意 $A,B\in\ell(\cP)$ 有 $B\in\cG_A$，即 $A\cap B\in\ell(\cP)$。所以 $\ell(\cP)$ 是 $\pi$-类。
\end{proof}

这是一个值得学习的证明。我们想证明的是一个关于两个变量 $A,B$ 的性质（$A\cap B\in\ell(\cP)$），直接证明做不到，于是分两次放大：(b) 先让 $A$ 留在小的集合族 $\cP$ 里，把 $B$ 放大到 $\ell(\cP)$；(c) 再利用交集的对称性，把 $A$ 也放大到 $\ell(\cP)$。每一次放大用的都是同一个事实，即“好集”构成 $\lambda$-类，而 $\ell(\cP)$ 是最小的 $\lambda$-类。以后证明乘积测度的性质时，我们还会见到同样的两次放大。

\begin{theorem}[Dynkin 的 $\pi$-$\lambda$ 定理]\label{thm:set-systems:pi-lambda}
设 $\cP$ 是 $\pi$-类，$\cL$ 是 $\lambda$-类，且 $\cP\subseteq\cL$。则 $\sigma(\cP)\subseteq\cL$。进一步，$\ell(\cP)=\sigma(\cP)$。
\end{theorem}
\begin{proof}
由引理~\ref{lem:set-systems:ell-is-pi}，$\ell(\cP)$ 是 $\pi$-类；它当然也是 $\lambda$-类。由引理~\ref{lem:set-systems:pi-lambda-sigma}，$\ell(\cP)$ 是 $\sigma$-代数。它包含 $\cP$，所以由 $\sigma(\cP)$ 的最小性，$\sigma(\cP)\subseteq\ell(\cP)$。前面已经说明 $\ell(\cP)\subseteq\sigma(\cP)$，故 $\ell(\cP)=\sigma(\cP)$。最后，$\cL$ 是包含 $\cP$ 的 $\lambda$-类，由 $\ell(\cP)$ 的最小性，$\sigma(\cP)=\ell(\cP)\subseteq\cL$。
\end{proof}

于是好集原理可以改写成下面的形式。

\begin{corollary}[$\pi$-$\lambda$ 形式的好集原理]\label{cor:set-systems:good-sets}
设 $\cP$ 是 $\pi$-类，$\cG\subseteq2^\Omega$。若 $\cP\subseteq\cG$ 且 $\cG$ 是 $\lambda$-类，则 $\sigma(\cP)\subseteq\cG$。
\end{corollary}
\begin{proof}
这就是定理~\ref{thm:set-systems:pi-lambda}。
\end{proof}

\begin{remark}[怎样使用 $\pi$-$\lambda$ 定理]
在应用中，我们总是按照同样的步骤来使用这个定理，同学们应该熟记：
\begin{enumerate}
\item 找一个生成所考虑的 $\sigma$-代数的 $\pi$-类 $\cP$，并在 $\cP$ 上验证所要的性质；
\item 令 $\cG$ 为具有这个性质的集合的全体，验证 $\cG$ 是 $\lambda$-类；
\item 由推论~\ref{cor:set-systems:good-sets} 得到 $\sigma(\cP)\subseteq\cG$。
\end{enumerate}
第 1 步通常是具体的计算，第 2 步通常是测度或积分的一般性质，第 3 步不需要做任何事情。下一讲的测度唯一性定理、第~\ref{lec:random-variables} 讲的单调类定理，以及以后关于独立性的判别，都是这样证明的。
\end{remark}

\begin{remark}
$\pi$-类的条件不能去掉。\egref{ex:set-systems:even} 中取 $\cC=\{\{1,2\},\{1,3\}\}$，则 $\cC\subseteq\cL$，$\cL$ 是 $\lambda$-类，但 $\sigma(\cC)$ 含有 $\{1\}=\{1,2\}\cap\{1,3\}$，从而 $\sigma(\cC)=2^\Omega\not\subseteq\cL$（$\{1\}$、$\{2\}=\{1,2\}\setminus\{1\}$、$\{3\}=\{1,3\}\setminus\{1\}$、$\{4\}=(\{1,2\}\cup\{1,3\})^c$ 都在 $\sigma(\cC)$ 中）。
\end{remark}

\begin{workedexample}\label{we:set-systems:pi-lambda-finite}
设 $\Omega=\{1,2,3,4\}$，$\mu,\nu$ 是 $\Omega$ 上的两个概率（即给每个点一个非负的权重，总和为 $1$，集合的概率是其中各点权重之和）。设 $\mu$ 与 $\nu$ 在 $\{1\},\{1,2\},\{1,2,3\}$ 上取值相同。用 $\pi$-$\lambda$ 定理证明 $\mu=\nu$，并直接验证这个结论。再说明：若只知道 $\mu,\nu$ 在 $\{1,2\},\{1,3\}$ 上相同，则推不出 $\mu=\nu$。
\end{workedexample}
\begin{solution}
\emph{用 $\pi$-$\lambda$ 定理。}按照上面“怎样使用 $\pi$-$\lambda$ 定理”的三步来做。
\begin{enumerate}
\item 取 $\cP=\{\{1\},\{1,2\},\{1,2,3\},\Omega\}$。这些集合是一条递增的链，任意两个的交是其中较小的一个，所以 $\cP$ 是 $\pi$-类。$\mu,\nu$ 在 $\cP$ 上相等（$\mu(\Omega)=\nu(\Omega)=1$）。$\cP$ 生成 $2^\Omega$：单点 $\{2\}=\{1,2\}\setminus\{1\}$，$\{3\}=\{1,2,3\}\setminus\{1,2\}$，$\{4\}=\Omega\setminus\{1,2,3\}$ 都在 $\sigma(\cP)$ 中，每个子集是单点的并。
\item 令 $\cD=\{A\subseteq\Omega:\mu(A)=\nu(A)\}$。本讲 $\lambda$-类一节中已经验证，$\cD$ 是 $\lambda$-类。在有限集上这只用到“不交集合的概率可加”。
\item 由定理~\ref{thm:set-systems:pi-lambda}，$2^\Omega=\sigma(\cP)\subseteq\cD$，即 $\mu=\nu$。
\end{enumerate}

\emph{直接验证。}记 $\mu$ 的权重为 $p_1,\dots,p_4$，$\nu$ 的为 $q_1,\dots,q_4$。由 $\mu(\{1\})=\nu(\{1\})$ 得 $p_1=q_1$；由 $\mu(\{1,2\})=\nu(\{1,2\})$ 得 $p_1+p_2=q_1+q_2$，从而 $p_2=q_2$；同理 $p_3=q_3$；最后由总和为 $1$ 得 $p_4=q_4$。这正是上面 $\{2\}=\{1,2\}\setminus\{1\}$ 等式子的数值版本。

\emph{反例。}在 $\lambda$-类一节开头的例子中，$\mu$ 是均匀分布，$\nu$ 给 $1,4$ 各 $\frac12$。两者在 $\{1,2\}$ 与 $\{1,3\}$ 上都等于 $\frac12$，但 $\mu(\{1\})=\frac14\neq\frac12=\nu(\{1\})$。原因是 $\{\{1,2\},\{1,3\},\Omega\}$ 不是 $\pi$-类：它的交 $\{1\}$ 不在其中。

这个例题是下一讲测度唯一性定理的缩影。“$\{1\},\{1,2\},\{1,2,3\}$”换成“一切 $(-\infty,x]$”，结论就是：$\R$ 上的概率由它的分布函数 $x\mapsto \P((-\infty,x])$ 唯一决定。
\end{solution}

\topic{本讲小结}

\begin{enumerate}
\item 概率空间 $(\Omega,\cF,\P)$ 中，$\cF$ 是“可以谈论概率的集合”的全体。它必须对补集和\emph{可数}并封闭（$\sigma$-代数），这样关于极限的事件（如 $\limsup_nA_n$、$\{S_n/n\to\frac12\}$）才是事件；又不能总取成全体子集。
\item 集合列的上极限 $\limsup_nA_n$ 是“$A_n$ 发生无穷多次”，下极限 $\liminf_nA_n$ 是“$A_n$ 从某项起总发生”。
\item 一般的 $\sigma$-代数通过“生成”得到：$\sigma(\cC)$ 是包含 $\cC$ 的最小 $\sigma$-代数。有限个集合生成的 $\sigma$-代数可以写出来（由划分的块组成）；一般情形写不出来，只能用好集原理。
\item $\R$ 上的 Borel $\sigma$-代数 $\cR$ 由开集生成，也由一切 $(-\infty,x]$ 生成。
\item $\pi$-$\lambda$ 定理：$\pi$-类 $\cP$ 包含于 $\lambda$-类 $\cL$，则 $\sigma(\cP)\subseteq\cL$。用法：在 $\pi$-类上验证，在 $\lambda$-类上推广。它在下一讲给出测度的唯一性（第~\ref{lec:measures} 讲），以后还给出单调类定理（第~\ref{lec:random-variables} 讲）与独立性的判别（第~\ref{lec:independence} 讲）。
\end{enumerate}

% =====================================================================
\homework{单调类定理，无穷 \texorpdfstring{$\sigma$}{σ}-代数的基数，可数生成}[单调类定理，无穷 \texorpdfstring{$\sigma$}{σ}-代数的基数，可数生成]
\label{hw:set-systems}


\exercisepart{基本习题}

\begin{exercises}{A}{课堂内容的补充}
\item\label{ex:set-systems:A1} 设 $A_n,B_n\subseteq\Omega$。
\begin{questions}
\item 证明 $\limsup_n(A_n\cup B_n)=\limsup_nA_n\cup\limsup_nB_n$。
\begin{solution}
用命题~\ref{prop:set-systems:limsup}(1)：$\omega\in\limsup_nC_n$ 当且仅当 $\omega$ 属于无穷多个 $C_n$。

若 $\omega$ 属于无穷多个 $A_n$，则它属于无穷多个 $A_n\cup B_n$（因为 $A_n\subseteq A_n\cup B_n$）；对 $B_n$ 同理。所以右边 $\subseteq$ 左边。

反之，设 $\omega$ 不属于右边，即 $\omega$ 只属于有限个 $A_n$，也只属于有限个 $B_n$。则存在 $N_1$ 使 $n\geqslant N_1$ 时 $\omega\notin A_n$，存在 $N_2$ 使 $n\geqslant N_2$ 时 $\omega\notin B_n$。于是 $n\geqslant\max\{N_1,N_2\}$ 时 $\omega\notin A_n\cup B_n$，$\omega$ 只属于有限个 $A_n\cup B_n$，不属于左边。所以左边 $\subseteq$ 右边。
\end{solution}
\item 证明 $\liminf_n(A_n\cap B_n)=\liminf_nA_n\cap\liminf_nB_n$。
\begin{solution}
用命题~\ref{prop:set-systems:limsup}(2)。若 $\omega$ 属于右边，存在 $N_1$ 使 $n\geqslant N_1$ 时 $\omega\in A_n$，存在 $N_2$ 使 $n\geqslant N_2$ 时 $\omega\in B_n$，于是 $n\geqslant\max\{N_1,N_2\}$ 时 $\omega\in A_n\cap B_n$，$\omega$ 属于左边。反之若从某个 $N$ 起 $\omega\in A_n\cap B_n$，则从 $N$ 起 $\omega\in A_n$，也从 $N$ 起 $\omega\in B_n$，$\omega$ 属于右边。

也可以对 A1-1) 取补集：由命题~\ref{prop:set-systems:limsup}(3)，$(\limsup_nC_n)^c=\liminf_nC_n^c$，把 A1-1) 中的 $A_n,B_n$ 换成 $A_n^c,B_n^c$，再两边取补，用 De Morgan 律即得。
\end{solution}
\item 证明 $\limsup_n(A_n\cap B_n)\subseteq\limsup_nA_n\cap\limsup_nB_n$，并举例说明等号可以不成立。
\begin{solution}
若 $\omega$ 属于无穷多个 $A_n\cap B_n$，则它属于这些 $A_n$，也属于这些 $B_n$，所以属于无穷多个 $A_n$ 与无穷多个 $B_n$。

反例：$n$ 为偶数时令 $A_n=\Omega$，$B_n=\varnothing$；$n$ 为奇数时令 $A_n=\varnothing$，$B_n=\Omega$。则 $A_n\cap B_n=\varnothing$ 对一切 $n$ 成立，左边是 $\varnothing$；而每个 $\omega$ 属于无穷多个 $A_n$（偶数项）和无穷多个 $B_n$（奇数项），右边是 $\Omega$。直观上，“$A$ 发生无穷多次”与“$B$ 发生无穷多次”不要求二者在\emph{同一时刻}发生。
\end{solution}
\end{questions}

\item\label{ex:set-systems:A2} 求下列集合列的上极限与下极限。
\begin{questions}
\item $\Omega=\R$，$A_n=[0,1]$（$n$ 为奇数），$A_n=[1,2]$（$n$ 为偶数）。
\begin{solution}
$x\in[0,1)$ 属于每个奇数项、不属于任何偶数项；$x\in(1,2]$ 属于每个偶数项、不属于任何奇数项；$x=1$ 属于每一项；其余的 $x$ 不属于任何一项。所以
\[
\limsup_nA_n=[0,2],\qquad\liminf_nA_n=\{1\}.
\]
\end{solution}
\item $A_n=(-\frac1n,\,1+(-1)^n]$。
\begin{solution}
$n$ 为偶数时 $A_n=(-\frac1n,2]$，$n$ 为奇数时 $A_n=(-\frac1n,0]$。逐点判断：
\begin{itemize}
\item $x=0$：属于每一项。
\item $0<x\leqslant2$：属于每个偶数项，不属于任何奇数项。所以属于上极限、不属于下极限。
\item $x<0$：$x\in A_n$ 要求 $-\frac1n<x$，即 $n<\frac1{\abs x}$，只有有限个 $n$。
\item $x>2$：不属于任何一项。
\end{itemize}
所以 $\limsup_nA_n=[0,2]$，$\liminf_nA_n=\{0\}$。
\end{solution}
\end{questions}

\item\label{ex:set-systems:A3} 证明：$\cA\subseteq2^\Omega$ 是代数，当且仅当 $\Omega\in\cA$，且 $A,B\in\cA$ 蕴含 $A\setminus B\in\cA$。
\begin{solution}
\emph{必要性。}若 $\cA$ 是代数，由命题~\ref{prop:set-systems:closure}(1)，$A\setminus B=A\cap B^c\in\cA$。

\emph{充分性。}设 $\Omega\in\cA$ 且 $\cA$ 对差集封闭。
\begin{itemize}
\item 补集：$A^c=\Omega\setminus A\in\cA$。
\item 交集：$A\cap B=A\setminus(A\setminus B)$。验证：$\omega$ 属于右边当且仅当 $\omega\in A$ 且 $\omega\notin A\setminus B$，即 $\omega\in A$ 且（$\omega\notin A$ 或 $\omega\in B$），即 $\omega\in A\cap B$。所以 $A\cap B\in\cA$。
\item 并集：$A\cup B=(A^c\cap B^c)^c$，由前两条属于 $\cA$。
\end{itemize}
所以 $\cA$ 是代数。
\end{solution}

\item\label{ex:set-systems:A4} 设 $\cF_1\subseteq\cF_2\subseteq\cdots$ 是 $\Omega$ 上的 $\sigma$-代数。
\begin{questions}
\item 证明 $\bigcup_n\cF_n$ 是代数。
\begin{solution}
记 $\cA=\bigcup_n\cF_n$。$\Omega\in\cF_1\subseteq\cA$。若 $A\in\cA$，则 $A\in\cF_n$ 对某个 $n$ 成立，于是 $A^c\in\cF_n\subseteq\cA$。若 $A,B\in\cA$，设 $A\in\cF_m$，$B\in\cF_n$，令 $k=\max\{m,n\}$，由递增性 $A,B\in\cF_k$，所以 $A\cup B\in\cF_k\subseteq\cA$。

注意递增性是用在最后一步的：两个集合来自不同的 $\cF_n$ 时，要找一个同时包含它们的 $\cF_k$。
\end{solution}
\item 设 $\Omega=\N$，$\cF_n=\sigma(\{1\},\{2\},\dots,\{n\})$。写出 $\cF_n$ 的全部元素，并证明 $\bigcup_n\cF_n$ 不是 $\sigma$-代数。
\begin{solution}
$\{1\},\dots,\{n\}$ 与 $R_n=\{n+1,n+2,\dots\}$ 构成 $\N$ 的划分，按例题（第~\pageref{we:set-systems:two-sets} 页）的方法，$\cF_n$ 由这 $n+1$ 块的并组成：
\[
\cF_n=\{S:\ S\subseteq\{1,\dots,n\}\}\cup\{S\cup R_n:\ S\subseteq\{1,\dots,n\}\},
\]
共 $2^{n+1}$ 个元素。$\cF_n\subseteq\cF_{n+1}$，因为 $R_n=\{n+1\}\cup R_{n+1}$。

$\bigcup_n\cF_n$ 中的集合要么是有限集，要么包含某个 $R_n$（从而补集有限）。偶数集 $E=\{2,4,6,\dots\}$ 既不有限，补集（奇数集）也不有限，所以 $E\notin\bigcup_n\cF_n$。但 $E=\bigcup_k\{2k\}$，每个 $\{2k\}\in\cF_{2k}\subseteq\bigcup_n\cF_n$。所以 $\bigcup_n\cF_n$ 对可数并不封闭，不是 $\sigma$-代数。（它恰好是\egref{ex:set-systems:finite-cofinite} 中的代数。）
\end{solution}
\end{questions}

\item\label{ex:set-systems:A5} 证明 $\cR^d$ 由下列每一个集合族生成。
\begin{questions}
\item 开矩形 $(a_1,b_1)\times\cdots\times(a_d,b_d)$ 的全体（提示：仿照引理~\ref{lem:set-systems:open-sets-R}，用有理顶点的开矩形）。
\begin{solution}
记开矩形全体为 $\cC_1$。开矩形是开集，所以 $\sigma(\cC_1)\subseteq\cR^d$。

反之，设 $U\subseteq\R^d$ 是开集，$x\in U$。取 $\varepsilon>0$ 使球 $\{y:\abs{y-x}<\varepsilon\}\subseteq U$。令 $\delta=\varepsilon/\sqrt d$，对每个 $i$ 取有理数 $p_i\in(x_i-\delta,x_i)$，$q_i\in(x_i,x_i+\delta)$。则 $x\in R=\prod_i(p_i,q_i)$；若 $y\in R$，则每个 $\abs{y_i-x_i}<\delta$，$\abs{y-x}<\sqrt{d\delta^2}=\varepsilon$，所以 $R\subseteq U$。于是 $U$ 是含于 $U$ 的有理顶点开矩形之并；有理顶点开矩形由 $2d$ 个有理数决定，只有可数个。所以 $U$ 是 $\cC_1$ 中可数个集合的并，属于 $\sigma(\cC_1)$，从而 $\cR^d\subseteq\sigma(\cC_1)$。
\end{solution}
\item 矩形 $(a_1,b_1]\times\cdots\times(a_d,b_d]$ 的全体。
\begin{solution}
记这样的矩形全体为 $\cC_2$，用向量记号 $(a,b]=\prod_i(a_i,b_i]$，$(a,b)=\prod_i(a_i,b_i)$，$\mathbf 1=(1,\dots,1)$。由
\[
(a,b]=\bigcap_{n=1}^\infty\Bigl(a,\,b+\tfrac1n\mathbf 1\Bigr),\qquad
(a,b)=\bigcup_{n>N_0}\Bigl(a,\,b-\tfrac1n\mathbf 1\Bigr]
\]
（$N_0$ 取得使 $\frac1{N_0}<\min_i(b_i-a_i)$）得 $\cC_2\subseteq\sigma(\cC_1)$，$\cC_1\subseteq\sigma(\cC_2)$。两个等式逐坐标验证，与命题~\ref{prop:set-systems:borel-generators} 证明中的一维情形相同：$y$ 属于左边当且仅当每个坐标满足相应的一维条件。由命题~\ref{prop:set-systems:generated}(3) 及 A5-1)，$\sigma(\cC_2)=\sigma(\cC_1)=\cR^d$。
\end{solution}
\item $\{(-\infty,x_1]\times\cdots\times(-\infty,x_d]: x_1,\dots,x_d\in\Q\}$。
\begin{solution}
记这个族为 $\cC_4$，记 $(-\infty,x]=\prod_i(-\infty,x_i]$。它是闭集，所以 $\sigma(\cC_4)\subseteq\cR^d$。只需证 $\cC_2\subseteq\sigma(\cC_4)$。

\emph{第一步：$x$ 为实向量时 $(-\infty,x]\in\sigma(\cC_4)$。}对每个 $i$ 取有理数列 $q^{(n)}_i\downarrow x_i$，$q^{(n)}_i>x_i$（构造同命题~\ref{prop:set-systems:borel-generators} 的证明），则 $(-\infty,x]=\bigcap_n(-\infty,q^{(n)}]$。

\emph{第二步：用“角”表示矩形。}设 $a<b$（逐坐标）。对 $i=1,\dots,d$，令 $b^{(i)}$ 为把 $b$ 的第 $i$ 个坐标换成 $a_i$ 所得的向量。我们断言
\[
(a,b]=(-\infty,b]\setminus\bigcup_{i=1}^d(-\infty,b^{(i)}].
\]
事实上，$y\in(-\infty,b^{(i)}]$ 当且仅当 $y\leqslant b$ 且 $y_i\leqslant a_i$。所以 $y$ 属于右边当且仅当 $y\leqslant b$，且对每个 $i$ 都有 $y_i>a_i$，这正是 $y\in(a,b]$。由第一步，右边属于 $\sigma(\cC_4)$。
\end{solution}
\item 由 A5-3) 说明 $\cR^d$ 是\term{可数生成的}（countably generated），即存在可数集合族 $\cC$ 使 $\sigma(\cC)=\cR^d$。
\begin{solution}
$\cC_4$ 中的集合由 $x\in\Q^d$ 决定，$\Q^d$ 可数，所以 $\cC_4$ 是可数族，而 A5-3) 说 $\sigma(\cC_4)=\cR^d$。
\end{solution}
\end{questions}

\item\label{ex:set-systems:A6} 设 $A\subseteq\N$。若极限 $\lim_{n\to\infty}\abs{A\cap\{1,\dots,n\}}/n$ 存在，称 $A$ 有\term{渐近密度}（asymptotic density）。令 $\cA$ 为有渐近密度的集合的全体。
\begin{questions}
\item $\cA$ 是否对补集封闭？是否对有限个不交集合的并封闭？
\begin{solution}
记 $[n]=\{1,\dots,n\}$，$d_n(A)=\abs{A\cap[n]}/n$。

补集：$\abs{A^c\cap[n]}=n-\abs{A\cap[n]}$，所以 $d_n(A^c)=1-d_n(A)$。若 $d_n(A)$ 收敛，$d_n(A^c)$ 也收敛。封闭。

不交并：若 $A\cap B=\varnothing$，则 $\abs{(A\cup B)\cap[n]}=\abs{A\cap[n]}+\abs{B\cap[n]}$，$d_n(A\cup B)=d_n(A)+d_n(B)$，两个收敛数列之和收敛。封闭。归纳可得对有限个两两不交集合的并封闭。
\end{solution}
\item $\cA$ 是否对可数个不交集合的并封闭？是否是 $\sigma$-代数？
\begin{solution}
不是。每个单点集 $\{k\}$ 有密度 $0$，而 $\N$ 的任何子集都是可数个单点集的不交并。所以只要存在一个没有渐近密度的集合，$\cA$ 就对可数不交并不封闭。下一小问中的 $E\cap B$ 就是这样的集合（更简单的例子是 $\bigcup_k[4^k,2\cdot4^k)$，计算与下一小问相同）。所以 $\cA$ 不是 $\sigma$-代数。
\end{solution}
\item $\cA$ 是否是代数？（提示：考虑偶数集与一个“在偶数和奇数之间按越来越长的段落交替选取”的集合的交。）
\begin{solution}
不是。我们找两个有密度的集合，它们的交没有密度。

令 $I_j=\{2^j,2^j+1,\dots,2^{j+1}-1\}$（$j\geqslant0$），这些段两两不交，并为 $\N$；$j\geqslant1$ 时 $I_j$ 有 $2^j$ 个数，从偶数开始，奇偶各半。令 $E$ 为偶数集，
\[
B=\bigcup_{j\ \text{偶}}(I_j\cap E)\ \cup\bigcup_{j\ \text{奇}}(I_j\cap E^c),
\]
即在第偶数段取偶数，在第奇数段取奇数。

\emph{$E$ 有密度 $\frac12$}：$\abs{E\cap[n]}=\lfloor n/2\rfloor$。

\emph{$B$ 有密度 $\frac12$。}设 $2^m\leqslant n<2^{m+1}$。$I_0=\{1\}$ 不含 $B$ 的点；$1\leqslant j\leqslant m-1$ 时 $B$ 在 $I_j$ 中恰有 $2^{j-1}$ 个点，共 $2^{m-1}-1$ 个；在 $I_m\cap[n]=\{2^m,\dots,n\}$ 中，$B$ 的点是其中的偶数或奇数，个数与 $\frac{n-2^m+1}2$ 相差不超过 $1$。所以 $\abs{\abs{B\cap[n]}-\frac n2}\leqslant2$，$d_n(B)\to\frac12$。

\emph{$E\cap B$ 没有密度。}$E\cap B=\bigcup_{j\ \text{偶},\,j\geqslant2}(I_j\cap E)$，第 $2i$ 段贡献 $2^{2i-1}$ 个点。取 $n_k=2^{2k+1}-1$（第 $2k$ 段的末尾），则
\[
\abs{E\cap B\cap[n_k]}=\sum_{i=1}^k2^{2i-1}=\frac{4^{k+1}-4}6,\qquad d_{n_k}(E\cap B)=\frac{(4^{k+1}-4)/6}{2\cdot4^k-1}\to\frac13 .
\]
取 $m_k=2^{2k+2}-1$（第 $2k+1$ 段的末尾），这一段不贡献新的点，所以
\[
d_{m_k}(E\cap B)=\frac{(4^{k+1}-4)/6}{4^{k+1}-1}\to\frac16 .
\]
两个子列的极限不同，$E\cap B$ 没有渐近密度。所以 $\cA$ 对交集不封闭；又由于 $E\cap B=(E^c\cup B^c)^c$ 而 $\cA$ 对补集封闭，$\cA$ 对并集也不封闭，不是代数。

这个例子说明，“渐近密度”虽然很像概率（有限可加、全集为 $1$），却不是定义在一个代数上的。
\end{solution}
\end{questions}

\item\label{ex:set-systems:A7} 设 $f:\Omega\to S$ 是映射，$\cS$ 是 $S$ 上的 $\sigma$-代数，$\cF$ 是 $\Omega$ 上的 $\sigma$-代数。
\begin{questions}
\item 证明 $f^{-1}(\cS)=\{f^{-1}(B): B\in\cS\}$ 是 $\Omega$ 上的 $\sigma$-代数。
\begin{solution}
原像保持一切集合运算：$f^{-1}(S)=\Omega$，$f^{-1}(B^c)=(f^{-1}(B))^c$，$f^{-1}(\bigcup_nB_n)=\bigcup_nf^{-1}(B_n)$。比如第二式：$\omega\in f^{-1}(B^c)\iff f(\omega)\notin B\iff\omega\notin f^{-1}(B)$。于是
\begin{itemize}
\item $\Omega=f^{-1}(S)\in f^{-1}(\cS)$；
\item 若 $A=f^{-1}(B)$，$B\in\cS$，则 $A^c=f^{-1}(B^c)$，$B^c\in\cS$；
\item 若 $A_n=f^{-1}(B_n)$，则 $\bigcup_nA_n=f^{-1}(\bigcup_nB_n)$，$\bigcup_nB_n\in\cS$。
\end{itemize}
\end{solution}
\item 证明 $\{B\subseteq S: f^{-1}(B)\in\cF\}$ 是 $S$ 上的 $\sigma$-代数。
\begin{solution}
记这个族为 $\cG$。$f^{-1}(S)=\Omega\in\cF$，所以 $S\in\cG$。若 $B\in\cG$，$f^{-1}(B^c)=(f^{-1}(B))^c\in\cF$，所以 $B^c\in\cG$。若 $B_n\in\cG$，$f^{-1}(\bigcup_nB_n)=\bigcup_nf^{-1}(B_n)\in\cF$，所以 $\bigcup_nB_n\in\cG$。
\end{solution}
\item 对任意 $\cC\subseteq2^S$，证明 $\sigma(f^{-1}(\cC))=f^{-1}(\sigma(\cC))$（提示：用 A7-2) 证明一个方向）。
\begin{solution}
“$\subseteq$”：由 A7-1)（把 $\cS$ 换成 $\sigma(\cC)$），$f^{-1}(\sigma(\cC))$ 是 $\sigma$-代数，它包含 $f^{-1}(\cC)$，由最小性得 $\sigma(f^{-1}(\cC))\subseteq f^{-1}(\sigma(\cC))$。

“$\supseteq$”：记 $\cF'=\sigma(f^{-1}(\cC))$。由 A7-2)（把 $\cF$ 换成 $\cF'$），$\cG=\{B\subseteq S: f^{-1}(B)\in\cF'\}$ 是 $S$ 上的 $\sigma$-代数。对 $C\in\cC$，$f^{-1}(C)\in\cF'$，所以 $\cC\subseteq\cG$，由最小性 $\sigma(\cC)\subseteq\cG$。这就是说，对每个 $B\in\sigma(\cC)$，$f^{-1}(B)\in\cF'$，即 $f^{-1}(\sigma(\cC))\subseteq\cF'$。

这里的“$\supseteq$”方向用的正是好集原理：好集是“原像落在 $\cF'$ 中的集合”。第~\ref{lec:random-variables} 讲判断随机变量可测性时，用的就是这个结论。
\end{solution}
\end{questions}

\item\label{ex:set-systems:A8} 设 $\Omega_0\subseteq\Omega$，$\cF$ 是 $\Omega$ 上的 $\sigma$-代数。
\begin{questions}
\item 证明 $\cF\cap\Omega_0=\{A\cap\Omega_0: A\in\cF\}$ 是 $\Omega_0$ 上的 $\sigma$-代数（称为\term{迹 $\sigma$-代数}，trace $\sigma$-algebra）。
\begin{solution}
注意补集要在 $\Omega_0$ 中取。$\Omega_0=\Omega\cap\Omega_0$。$\Omega_0\setminus(A\cap\Omega_0)=A^c\cap\Omega_0$，$A^c\in\cF$。$\bigcup_n(A_n\cap\Omega_0)=(\bigcup_nA_n)\cap\Omega_0$。三条都满足。

也可以用 A7-1)：令 $\iota:\Omega_0\to\Omega$ 为包含映射 $\iota(\omega)=\omega$，则 $\iota^{-1}(A)=A\cap\Omega_0$，$\cF\cap\Omega_0=\iota^{-1}(\cF)$。
\end{solution}
\item 对 $\cC\subseteq2^\Omega$，证明 $\sigma_{\Omega_0}(\cC\cap\Omega_0)=\sigma(\cC)\cap\Omega_0$，其中左边表示在 $\Omega_0$ 上生成的 $\sigma$-代数。
\begin{solution}
仍用包含映射 $\iota$，$\iota^{-1}(C)=C\cap\Omega_0$。由 A7-3)（$f=\iota$，$S=\Omega$），
\[
\sigma_{\Omega_0}(\cC\cap\Omega_0)=\sigma_{\Omega_0}(\iota^{-1}(\cC))=\iota^{-1}(\sigma(\cC))=\sigma(\cC)\cap\Omega_0 .
\]
\end{solution}
\end{questions}

\item\label{ex:set-systems:A9} 设 $\cL$ 是 $\lambda$-类。
\begin{questions}
\item 证明：若 $A,B\in\cL$ 且 $A\cap B=\varnothing$，则 $A\cup B\in\cL$。
\begin{solution}
$B^c=\Omega\setminus B\in\cL$。$A\cap B=\varnothing$ 意味着 $A\subseteq B^c$，所以 $B^c\setminus A\in\cL$，再取补：$(B^c\setminus A)^c=(B^c\cap A^c)^c=B\cup A\in\cL$。（这就是命题~\ref{prop:set-systems:lambda-equivalent} 证明中的一步。）
\end{solution}
\item 举例说明 $A,B\in\cL$ 一般推不出 $A\cup B\in\cL$。
\begin{solution}
取\egref{ex:set-systems:even} 中 $\{1,2,3,4\}$ 的偶数元子集组成的 $\lambda$-类。$\{1,2\},\{1,3\}\in\cL$，但 $\{1,2\}\cup\{1,3\}=\{1,2,3\}$ 有 $3$ 个元素，不属于 $\cL$。
\end{solution}
\end{questions}
\end{exercises}

\begin{probleme}{B}{单调类定理}
$\pi$-$\lambda$ 定理有一个历史上更早的“孪生”版本。称 $\cM\subseteq2^\Omega$ 为\term{单调类}（monotone class），如果 $A_n\in\cM$，$A_n\uparrow A$ 蕴含 $A\in\cM$，并且 $A_n\in\cM$，$A_n\downarrow A$ 蕴含 $A\in\cM$。本题证明：\emph{若 $\cA$ 是代数，则包含 $\cA$ 的最小单调类 $m(\cA)$ 等于 $\sigma(\cA)$。}
\begin{questions}
\item 证明任意一族单调类的交是单调类，从而 $m(\cA)$ 有定义，并且 $m(\cA)\subseteq\sigma(\cA)$。
\begin{solution}
设 $\cM_i$（$i\in I$）都是单调类，$A_n\in\bigcap_i\cM_i$，$A_n\uparrow A$。对每个 $i$，$A_n\in\cM_i$，所以 $A\in\cM_i$，从而 $A\in\bigcap_i\cM_i$；递减的情形相同。$2^\Omega$ 是包含 $\cA$ 的单调类，所以 $m(\cA)=\bigcap\{\cM:\cM\ \text{是单调类},\ \cA\subseteq\cM\}$ 有定义，是包含 $\cA$ 的最小单调类。$\sigma(\cA)$ 对可数并、可数交封闭，当然是单调类，并包含 $\cA$，所以 $m(\cA)\subseteq\sigma(\cA)$。
\end{solution}
\item 证明：若一个代数同时是单调类，则它是 $\sigma$-代数。
\begin{solution}
设 $\cG$ 是代数也是单调类，$A_1,A_2,\dots\in\cG$。令 $B_n=A_1\cup\cdots\cup A_n$，由代数对有限并封闭，$B_n\in\cG$；又 $B_n\uparrow\bigcup_kA_k$，由单调类的性质，$\bigcup_kA_k\in\cG$。其余两条是代数的定义。
\end{solution}
\item 对 $A\subseteq\Omega$，令 $\cM_A=\{B\in m(\cA): A\cap B,\ A\setminus B,\ B\setminus A\in m(\cA)\}$。证明：若 $A\in m(\cA)$，则 $\cM_A$ 是单调类。
\begin{solution}
设 $B_n\in\cM_A$，$B_n\uparrow B$。则 $B\in m(\cA)$（$m(\cA)$ 是单调类）。又
\[
A\cap B_n\uparrow A\cap B,\qquad A\setminus B_n\downarrow A\setminus B,\qquad B_n\setminus A\uparrow B\setminus A,
\]
左边三列都在 $m(\cA)$ 中，由单调性，右边三个集合都在 $m(\cA)$ 中，所以 $B\in\cM_A$。验证第二式：$A\setminus B_n=A\cap B_n^c$，$B_n^c$ 递减，交为 $B^c$。递减的情形 $B_n\downarrow B$ 完全对称：$A\cap B_n\downarrow A\cap B$，$A\setminus B_n\uparrow A\setminus B$，$B_n\setminus A\downarrow B\setminus A$。

（这个论证其实对任何 $A$ 都成立；$A\in m(\cA)$ 的假设在 B5) 中才用到。）
\end{solution}
\item 证明：若 $A\in\cA$，则 $\cA\subseteq\cM_A$，从而 $m(\cA)\subseteq\cM_A$。
\begin{solution}
任取 $B\in\cA$。$\cA$ 是代数，所以 $A\cap B$，$A\setminus B$，$B\setminus A$ 都在 $\cA\subseteq m(\cA)$ 中，又 $B\in m(\cA)$，所以 $B\in\cM_A$。于是 $\cA\subseteq\cM_A$。由 B3)，$\cM_A$ 是单调类，由 $m(\cA)$ 的最小性得 $m(\cA)\subseteq\cM_A$。
\end{solution}
\item 利用 $B\in\cM_A\iff A\in\cM_B$，证明：若 $A\in m(\cA)$，则 $m(\cA)\subseteq\cM_A$。由此说明 $m(\cA)$ 对交集和差集封闭。
\begin{solution}
先说明对称性：当 $A,B\in m(\cA)$ 时，$B\in\cM_A$ 的条件是“$A\cap B,A\setminus B,B\setminus A\in m(\cA)$”，$A\in\cM_B$ 的条件是“$B\cap A,B\setminus A,A\setminus B\in m(\cA)$”，两者相同。

设 $A\in m(\cA)$。任取 $B\in\cA$，由 B4)（把 B4) 中的 $A$ 换成 $B$），$m(\cA)\subseteq\cM_B$，所以 $A\in\cM_B$，由对称性 $B\in\cM_A$。这说明 $\cA\subseteq\cM_A$。由 B3) 与最小性，$m(\cA)\subseteq\cM_A$。

于是对任意 $A,B\in m(\cA)$，$B\in\cM_A$，即 $A\cap B$，$A\setminus B\in m(\cA)$。这和引理~\ref{lem:set-systems:ell-is-pi} 的证明一样，是“两次放大”。
\end{solution}
\item 证明 $\Omega\in m(\cA)$，从而 $m(\cA)$ 是代数。结合 B2) 得到 $m(\cA)=\sigma(\cA)$。
\begin{solution}
$\Omega\in\cA\subseteq m(\cA)$。由 B5)，$m(\cA)$ 对差集封闭，由习题~\exref{ex:set-systems:A3}，$m(\cA)$ 是代数。由 B2)，它是 $\sigma$-代数，包含 $\cA$，所以 $\sigma(\cA)\subseteq m(\cA)$。与 B1) 合起来得 $m(\cA)=\sigma(\cA)$。
\end{solution}
\end{questions}
最后比较两个定理：
\begin{questions}[resume]
\item 用单调类定理证明：若 $\cA$ 是代数，$\cL$ 是包含 $\cA$ 的 $\lambda$-类，则 $\sigma(\cA)\subseteq\cL$。（提示：$\lambda$-类对递增并封闭；它对递减交也封闭吗？）
\begin{solution}
$\lambda$-类对递增并封闭。设 $A_n\in\cL$，$A_n\downarrow A$。则 $A_n^c=\Omega\setminus A_n\in\cL$，且 $A_n^c\uparrow A^c$，所以 $A^c\in\cL$，再取补 $A\in\cL$。于是 $\cL$ 是单调类，包含 $\cA$，所以 $\sigma(\cA)=m(\cA)\subseteq\cL$。
\end{solution}
\item 举例说明：若 $\cA$ 只是 $\pi$-类而不是代数，则 $m(\cA)$ 可以严格小于 $\sigma(\cA)$。（提示：在 $\Omega=\{1,2\}$ 上取 $\cA=\{\{1\}\}$。）
\begin{solution}
$\cA=\{\{1\}\}$ 是 $\pi$-类。$\cA$ 本身是单调类：其中的单调列只能是常数列 $\{1\},\{1\},\dots$，极限仍是 $\{1\}$。所以 $m(\cA)=\{\{1\}\}$。而 $\sigma(\cA)=\{\varnothing,\{1\},\{2\},\Omega\}$。所以 $m(\cA)\subsetneq\sigma(\cA)$。

比较：$\pi$-$\lambda$ 定理只要求生成元是 $\pi$-类，这比代数容易验证（比如 $\{(-\infty,x]\}$ 是 $\pi$-类，但不是代数），所以在概率论中更常用。
\end{solution}
\end{questions}
\end{probleme}

\begin{probleme}{C}{无穷 $\sigma$-代数不可数}
本题证明：\emph{若 $\sigma$-代数 $\cF$ 有无穷多个元素，则 $\cF$ 的基数不小于连续统的基数；特别地，不存在可数无穷的 $\sigma$-代数。}
\begin{questions}
\item 对 $\omega\in\Omega$，令 $[\omega]=\bigcap\{A\in\cF:\omega\in A\}$。说明 $[\omega]$ 不一定属于 $\cF$。现假设 $\cF$ 是可数的，证明 $[\omega]\in\cF$。
\begin{solution}
\emph{反例。}取不可数集 $\Omega$（如 $\R$），固定一点 $\omega_0$，令
\[
\cF=\{A:\ A\ \text{可数且}\ \omega_0\notin A\}\cup\{A:\ A^c\ \text{可数且}\ \omega_0\in A\}.
\]
它是 $\sigma$-代数：$\Omega$ 属于第二类；取补把两类互换；一列集合若都属于第一类，并集可数且不含 $\omega_0$；若有一个属于第二类，并集含 $\omega_0$ 且补集可数。含 $\omega_0$ 的集合都是第二类的，它们的交是 $\{\omega_0\}$（对每个 $\omega\neq\omega_0$，$\Omega\setminus\{\omega\}$ 是含 $\omega_0$ 的第二类集合）。所以 $[\omega_0]=\{\omega_0\}$，它可数却含 $\omega_0$，不属于 $\cF$。

\emph{$\cF$ 可数时。}此时 $\{A\in\cF:\omega\in A\}$ 是可数族，$[\omega]$ 是 $\cF$ 中可数个集合的交，属于 $\cF$。
\end{solution}
\item 仍假设 $\cF$ 可数。证明：对任意 $\omega,\omega'$，$[\omega]$ 与 $[\omega']$ 或者相等，或者不交。（提示：若 $\omega'\notin[\omega]$，则存在 $A\in\cF$ 使 $\omega\in A$，$\omega'\notin A$，考虑 $A^c$。）
\begin{solution}
分两种情形。

\emph{情形一：$\omega'\notin[\omega]$。}按定义存在 $A\in\cF$ 使 $\omega\in A$，$\omega'\notin A$。于是 $[\omega]\subseteq A$；又 $\omega'\in A^c\in\cF$，所以 $[\omega']\subseteq A^c$。两者不交。

\emph{情形二：$\omega'\in[\omega]$。}这意味着每个含 $\omega$ 的 $A\in\cF$ 都含 $\omega'$，所以 $\{A:\omega\in A\}\subseteq\{A:\omega'\in A\}$，对更大的族取交得到更小的集合：$[\omega']\subseteq[\omega]$。再证 $\omega\in[\omega']$：如若不然，由情形一（交换 $\omega,\omega'$ 的角色）存在 $A\in\cF$ 使 $\omega'\in A$，$\omega\notin A$，于是 $\omega\in A^c$，从而 $\omega'\in A^c$（因为 $\omega'\in[\omega]\subseteq A^c$），与 $\omega'\in A$ 矛盾。所以 $\omega\in[\omega']$，同理得 $[\omega]\subseteq[\omega']$。两者相等。

（这一小问其实不需要 $\cF$ 可数。）
\end{solution}
\item 仍假设 $\cF$ 可数。设 $\{[\omega]:\omega\in\Omega\}$ 中互不相同的集合为 $E_1,E_2,\dots$（有限或可数个）。证明每个 $A\in\cF$ 都是若干个 $E_i$ 的并。
\begin{solution}
设 $A\in\cF$。对每个 $\omega\in A$，$A$ 是含 $\omega$ 的 $\cF$ 中集合之一，所以 $[\omega]\subseteq A$；又 $\omega\in[\omega]$。于是 $A=\bigcup_{\omega\in A}[\omega]$，右边是若干个 $E_i$ 的并。（$E_i$ 至多可数个，是因为它们是 $\cF$ 中两两不交的非空集合，而 $\cF$ 可数。）
\end{solution}
\item 由 C3) 和\egref{ex:set-systems:partition} 证明：若 $\cF$ 可数，则 $\cF$ 是有限的。
\begin{solution}
由 C1)–C3)，$E_1,E_2,\dots$ 是 $\Omega$ 的划分，每块属于 $\cF$。若块有无穷多个，则对每个 $I\subseteq\N$，$\bigcup_{i\in I}E_i$ 是可数并，属于 $\cF$；不同的 $I$ 给出不同的集合（各块非空且不交）。$\N$ 的子集有不可数多个，与 $\cF$ 可数矛盾。所以只有有限个块 $E_1,\dots,E_k$，由 C3)，$\cF$ 中每个集合是若干块的并，$\cF$ 至多有 $2^k$ 个元素。
\end{solution}
\item 现去掉可数的假设，设 $\cF$ 是无穷 $\sigma$-代数。证明存在一列两两不交的非空集合 $B_1,B_2,\dots\in\cF$。（提示：若 $\cF$ 无穷，则存在 $A\in\cF$ 使 $\cF\cap A=\{C\cap A: C\in\cF\}$ 和 $\cF\cap A^c$ 中至少一个是无穷的，而且 $A$、$A^c$ 都非空；归纳地往下取。）
\begin{solution}
\emph{第一步：一个切分引理。}设 $\Omega'\in\cF$，且 $\cF\cap\Omega'$ 是无穷的。我们断言：存在 $A\in\cF$，$A\subseteq\Omega'$，使 $A$ 与 $\Omega'\setminus A$ 都非空，并且 $\cF\cap A$ 与 $\cF\cap(\Omega'\setminus A)$ 中至少有一个是无穷的。

由于 $\cF\cap\Omega'$ 无穷，其中有一个集合 $A$ 既不是 $\varnothing$ 也不是 $\Omega'$；因为 $\Omega'\in\cF$，$A=C\cap\Omega'\in\cF$。映射 $D\mapsto(D\cap A,\ D\cap(\Omega'\setminus A))$ 把 $\cF\cap\Omega'$ 单射地映入 $(\cF\cap A)\times(\cF\cap(\Omega'\setminus A))$（$D$ 是两部分的并）。若后两个都有限，$\cF\cap\Omega'$ 也有限，矛盾。

\emph{第二步：归纳。}令 $\Omega_0=\Omega$。设已取得 $\Omega_k\in\cF$，$\cF\cap\Omega_k$ 无穷。用第一步得到 $A\subseteq\Omega_k$，在 $A$ 与 $\Omega_k\setminus A$ 中，取迹为无穷的那一个（都无穷时任取一个）记为 $\Omega_{k+1}$，另一个记为 $B_{k+1}$。则 $B_{k+1}\in\cF$ 非空，$B_{k+1}\subseteq\Omega_k\setminus\Omega_{k+1}$。

\emph{第三步：不交性。}若 $k<l$，则 $B_l\subseteq\Omega_{l-1}\subseteq\Omega_{k+1}$，而 $B_{k+1}\cap\Omega_{k+1}=\varnothing$，所以 $B_{k+1}\cap B_l=\varnothing$。
\end{solution}
\item 利用 C5) 中的 $B_n$，构造从 $2^\N$ 到 $\cF$ 的单射，从而完成证明。
\begin{solution}
令 $\Phi(I)=\bigcup_{n\in I}B_n$（$I\subseteq\N$）。可数并属于 $\cF$。若 $I\neq J$，比如 $n\in I\setminus J$，则 $B_n\subseteq\Phi(I)$，而 $B_n$ 非空且与 $\Phi(J)$ 不交，所以 $\Phi(I)\neq\Phi(J)$。于是 $\Phi$ 是单射，$\abs\cF\geqslant\abs{2^\N}=\abs\R$。
\end{solution}
\end{questions}
\end{probleme}

\begin{probleme}{D}{生成的 $\sigma$-代数中的集合只依赖可数个生成元}
设 $\cC\subseteq2^\Omega$ 是任意集合族（可以不可数）。本题证明：\emph{对每个 $A\in\sigma(\cC)$，存在可数子族 $\cC_0\subseteq\cC$ 使 $A\in\sigma(\cC_0)$。}这一结论在研究 $[0,1]$ 上的连续函数空间、无穷乘积空间时十分有用。
\begin{questions}
\item 令 $\cG=\bigcup\{\sigma(\cC_0):\ \cC_0\subseteq\cC\ \text{可数}\}$。证明 $\cC\subseteq\cG\subseteq\sigma(\cC)$。
\begin{solution}
对 $C\in\cC$，取 $\cC_0=\{C\}$，则 $C\in\sigma(\cC_0)\subseteq\cG$。对可数的 $\cC_0\subseteq\cC$，由命题~\ref{prop:set-systems:generated}(2)，$\sigma(\cC_0)\subseteq\sigma(\cC)$，所以 $\cG\subseteq\sigma(\cC)$。
\end{solution}
\item 证明 $\Omega\in\cG$，且 $\cG$ 对补集封闭。
\begin{solution}
$\Omega\in\sigma(\varnothing)=\{\varnothing,\Omega\}$，空族是可数的。若 $A\in\cG$，则 $A\in\sigma(\cC_0)$ 对某个可数 $\cC_0$ 成立，于是 $A^c\in\sigma(\cC_0)\subseteq\cG$。
\end{solution}
\item 设 $A_n\in\cG$，$A_n\in\sigma(\cC_n)$，$\cC_n$ 可数。说明 $\bigcup_n\cC_n$ 可数，并证明 $\bigcup_nA_n\in\sigma(\bigcup_n\cC_n)$。
\begin{solution}
可数个可数集的并可数，所以 $\cC'=\bigcup_n\cC_n$ 是 $\cC$ 的可数子族。由 $\cC_n\subseteq\cC'$ 得 $\sigma(\cC_n)\subseteq\sigma(\cC')$，所以每个 $A_n\in\sigma(\cC')$，它们的可数并也属于 $\sigma(\cC')\subseteq\cG$。
\end{solution}
\item 由 D1)–D3) 得到 $\cG=\sigma(\cC)$。
\begin{solution}
D2)、D3) 说明 $\cG$ 是 $\sigma$-代数，D1) 说明它包含 $\cC$，由最小性 $\sigma(\cC)\subseteq\cG$；反包含是 D1)。
\end{solution}
\end{questions}
作为应用，设 $T$ 是不可数集，$\Omega=\R^T$ 为一切函数 $f:T\to\R$ 的集合，$\cC=\{\{f: f(t)\leqslant x\}: t\in T,\ x\in\R\}$。
\begin{questions}[resume]
\item 证明：对每个 $A\in\sigma(\cC)$，存在可数集 $T_0\subseteq T$，使得若 $f,g\in\Omega$ 在 $T_0$ 上相等，则 $f\in A\iff g\in A$。
\begin{solution}
由 D4)，$A\in\sigma(\cC_0)$，$\cC_0=\{\{f:f(t_k)\leqslant x_k\}:k\in\N\}$ 可数。令 $T_0=\{t_k:k\in\N\}$。令
\[
\cH=\{B\subseteq\Omega:\ \text{若}\ f=g\ \text{在}\ T_0\ \text{上},\ \text{则}\ f\in B\iff g\in B\}.
\]
$\cH$ 是 $\sigma$-代数：$\Omega\in\cH$；条件“$f\in B\iff g\in B$”对 $B^c$ 与对 $B$ 等价；若对每个 $B_n$ 成立，则 $f\in\bigcup_nB_n\iff$ 存在 $n$ 使 $f\in B_n\iff$ 存在 $n$ 使 $g\in B_n\iff g\in\bigcup_nB_n$。$\cC_0\subseteq\cH$：$f\in\{f(t_k)\leqslant x_k\}$ 只依赖于 $f(t_k)$，$t_k\in T_0$。所以 $\sigma(\cC_0)\subseteq\cH$，$A\in\cH$。
\end{solution}
\item 取 $T=[0,1]$。证明 $\{f\in\R^{[0,1]}: f\ \text{在}\ [0,1]\ \text{上连续}\}\notin\sigma(\cC)$，也有 $\{f: \sup_{t\in[0,1]}f(t)\leqslant 1\}\notin\sigma(\cC)$。
\begin{solution}
设连续函数集 $A\in\sigma(\cC)$，取 D5) 中的可数集 $T_0$。$[0,1]$ 不可数，取 $s\in[0,1]\setminus T_0$。令 $f\equiv0$，$g=\one_{\{s\}}$。二者在 $T_0$ 上都为 $0$，所以 $f\in A\iff g\in A$。但 $f$ 连续，$g$ 在 $s$ 处不连续，矛盾。

对第二个集合，取同样的 $f\equiv0$ 与 $g=2\cdot\one_{\{s\}}$：$\sup f=0\leqslant1$，$\sup g=2>1$，同样矛盾。
\end{solution}
\end{questions}
这说明在一切函数构成的空间上，由“有限个时刻的取值”生成的 $\sigma$-代数太小，连“路径连续”都不是事件。下一讲会看到，把空间换成连续函数空间 $C[0,1]$，情况就好得多。
\end{probleme}

\thinkingpart

\begin{probleme}[问题]{E}{可数生成的 $\sigma$-代数与原子}
设 $(\Omega,\cF)$ 是可测空间。称 $\cF$ 可数生成，若存在可数族 $\cC=\{C_1,C_2,\dots\}$ 使 $\sigma(\cC)=\cF$。
\begin{questions}
\item 设 $\cF=\sigma(\cC)$，$\cC$ 如上。在 $\Omega$ 上定义 $\omega\sim\omega'$ 当且仅当对一切 $n$ 有 $\one_{C_n}(\omega)=\one_{C_n}(\omega')$。证明 $\sim$ 的每个等价类属于 $\cF$。
\begin{solution}
固定 $\omega$。对每个 $n$，令 $D_n=C_n$（若 $\omega\in C_n$），$D_n=C_n^c$（若 $\omega\notin C_n$）。则 $\omega'\sim\omega$ 当且仅当对每个 $n$，$\omega'$ 与 $\omega$ 同在或同不在 $C_n$ 中，即 $\omega'\in D_n$。所以 $\omega$ 的等价类等于 $\bigcap_nD_n$，是可数个 $\cF$ 中集合的交，属于 $\cF$。
\end{solution}
\item 证明：若 $A\in\cF$，$\omega\in A$，$\omega\sim\omega'$，则 $\omega'\in A$。（提示：好集原理。）于是每个等价类是 $\cF$ 的\term{原子}（atom）：它的可测子集只有空集和它自身。
\begin{solution}
令 $\cG=\{A\in\cF:\ \text{对一切}\ \omega\sim\omega',\ \omega\in A\Rightarrow\omega'\in A\}$。
\begin{itemize}
\item $C_n\in\cG$：由 $\sim$ 的定义，$\one_{C_n}(\omega)=\one_{C_n}(\omega')$。
\item $\Omega\in\cG$ 显然。
\item 补集：设 $A\in\cG$，$\omega\in A^c$，$\omega'\sim\omega$。若 $\omega'\in A$，由 $\sim$ 对称及 $A\in\cG$ 得 $\omega\in A$，矛盾。所以 $\omega'\in A^c$。
\item 可数并：若 $\omega\in\bigcup_nA_n$，则 $\omega\in A_n$ 对某个 $n$，从而 $\omega'\in A_n$。
\end{itemize}
所以 $\cG$ 是包含 $\cC$ 的 $\sigma$-代数，$\cF\subseteq\cG$。

原子：设 $E$ 是一个等价类，$B\in\cF$，$B\subseteq E$，$B\neq\varnothing$。取 $\omega\in B$，则 $E$ 中每个 $\omega'$ 都满足 $\omega'\sim\omega$，由刚证的结论 $\omega'\in B$。所以 $B=E$。
\end{solution}
\item 证明 $\R$ 上的可数－余可数 $\sigma$-代数不是可数生成的，尽管它比 $\cR$ “小”。（提示：若它由 $C_1,C_2,\dots$ 生成，每个 $C_n$ 可以取为可数集或余可数集，考察 E1) 中的等价类。）
\begin{solution}
设可数－余可数 $\sigma$-代数 $\cF$ 由 $C_1,C_2,\dots$ 生成，每个 $C_n\in\cF$ 可数或余可数。令 $D_n=C_n$（若 $C_n$ 可数），$D_n=C_n^c$（否则）。$\sigma(\{D_n\})=\sigma(\{C_n\})=\cF$（两族互相属于对方生成的 $\sigma$-代数）。$D=\bigcup_nD_n$ 可数，所以 $\R\setminus D$ 不可数。任取 $x\neq y$ 属于 $\R\setminus D$：对每个 $n$，$x,y\notin D_n$，所以 $x\sim y$（关于生成元 $D_n$）。由 E2)，$\{x\}\in\cF$ 是含 $x$ 的可测集，必须包含 $y$，矛盾。

这说明“可数生成”与“大小”无关：$\cR$ 比可数－余可数 $\sigma$-代数大，却是可数生成的（习题~\exref{ex:set-systems:A5}）。
\end{solution}
\end{questions}
\end{probleme}

\exercisepart{课外知识：Borel 集的基数}

\begin{probleme}{F}{$\cR$ 与 $\R$ 等势}
本题证明 $\abs{\cR}=\abs{\R}$，从而 $\cR\neq2^\R$（因为 $\abs{2^\R}>\abs{\R}$）。需要用到下面这个关于序数的事实，这里不加证明地承认：\emph{存在一个不可数的良序集 $(\Lambda,<)$，它的每个真前段 $\{\beta:\beta<\alpha\}$ 都是可数的；并且 $\Lambda$ 的每个可数子集都有严格上界。}（$\Lambda$ 就是第一个不可数序数 $\omega_1$。）
\begin{questions}
\item 令 $\cE_0$ 为 $\R$ 中端点为有理数的开区间的全体连同空集。对 $\alpha\in\Lambda$，若 $\alpha$ 是最小元，令 $\cE_\alpha=\cE_0$；否则令 $\cE_\alpha$ 为从 $\bigcup_{\beta<\alpha}\cE_\beta$ 中取可数多个集合作可数并、或者取一个集合作补集所得到的集合的全体。证明对一切 $\alpha$，$\cE_\alpha\subseteq\cR$。
\begin{solution}
用超限归纳法：若结论不对，因为 $\Lambda$ 良序，有一个最小的 $\alpha$ 使 $\cE_\alpha\not\subseteq\cR$。$\alpha$ 不是最小元，因为 $\cE_0$ 中的集合是开集。于是对一切 $\beta<\alpha$，$\cE_\beta\subseteq\cR$，而 $\cE_\alpha$ 中的集合是 $\cR$ 中集合的可数并或补集，仍属于 $\cR$，矛盾。
\end{solution}
\item 证明 $\cE=\bigcup_{\alpha\in\Lambda}\cE_\alpha$ 是 $\sigma$-代数。（提示：可数个 $\alpha_n$ 有严格上界。）从而 $\cE=\cR$。
\begin{solution}
\emph{补集。}设 $A\in\cE_\alpha$。单元素集 $\{\alpha\}$ 有严格上界 $\alpha'>\alpha$，于是 $A\in\bigcup_{\beta<\alpha'}\cE_\beta$，$A^c\in\cE_{\alpha'}\subseteq\cE$。

\emph{全集。}$\varnothing\in\cE_0$，由上一条 $\R=\varnothing^c\in\cE$。

\emph{可数并。}设 $A_n\in\cE_{\alpha_n}$。可数集 $\{\alpha_n\}$ 有严格上界 $\alpha$，于是每个 $A_n\in\bigcup_{\beta<\alpha}\cE_\beta$，$\bigcup_nA_n\in\cE_\alpha$。

所以 $\cE$ 是 $\sigma$-代数，它包含 $\cE_0$，而 $\sigma(\cE_0)=\cR$（引理~\ref{lem:set-systems:open-sets-R}：开集是有理端点开区间的可数并），所以 $\cR\subseteq\cE$；由 F1)，$\cE\subseteq\cR$。
\end{solution}
\item 用超限归纳法证明对一切 $\alpha$，$\abs{\cE_\alpha}\leqslant\abs{\R}$。（可以使用 $\abs{\R^\N}=\abs{\R}$ 以及“不超过 $\abs{\R}$ 个基数不超过 $\abs{\R}$ 的集合之并，基数不超过 $\abs{\R}$”。）
\begin{solution}
$\cE_0$ 可数。设对一切 $\beta<\alpha$ 有 $\abs{\cE_\beta}\leqslant\abs\R$。前段 $\{\beta<\alpha\}$ 可数，所以 $\cU=\bigcup_{\beta<\alpha}\cE_\beta$ 是可数个基数 $\leqslant\abs\R$ 的集合之并，$\abs\cU\leqslant\abs\R$。$\cE_\alpha$ 中的元素或者是 $\bigcup_nA_n$，$(A_n)\in\cU^\N$，或者是 $A^c$，$A\in\cU$，所以 $\cE_\alpha$ 是 $\cU^\N\cup\cU$ 的像，
\[
\abs{\cE_\alpha}\leqslant\abs{\cU^\N}+\abs\cU\leqslant\abs{\R^\N}+\abs\R=\abs\R .
\]
\end{solution}
\item 由 $\abs{\Lambda}\leqslant\abs{\R}$ 得到 $\abs{\cR}\leqslant\abs{\R}$，再证明反向不等式。
\begin{solution}
（$\abs\Lambda\leqslant\abs\R$：$\Lambda$ 是最小的不可数良序集的基数，而 $\R$ 不可数，这是集合论的标准事实，这里承认。）由 F2)、F3)，$\cR=\bigcup_{\alpha\in\Lambda}\cE_\alpha$ 是不超过 $\abs\R$ 个基数不超过 $\abs\R$ 的集合之并，$\abs\cR\leqslant\abs\R$。反向：$x\mapsto\{x\}$ 是从 $\R$ 到 $\cR$ 的单射。所以 $\abs\cR=\abs\R$。

由 Cantor 定理 $\abs{2^\R}>\abs\R$，所以 $\R$ 有不是 Borel 集的子集。注意这个论证只说明非 Borel 集存在，没有写出一个来；第~\ref{lec:extension-of-measures} 讲的作业会构造一个（借助选择公理）。
\end{solution}
\end{questions}
\end{probleme}
